\documentclass[journal]{IEEEtran}

% ── Packages ──
\usepackage[utf8]{inputenc}
\usepackage[T1]{fontenc}
\usepackage{amsmath,amssymb,amsfonts}
\usepackage{graphicx}
\usepackage{booktabs}
\usepackage{multirow}
\usepackage[dvipsnames]{xcolor}
\usepackage{hyperref}
\usepackage{cleveref}
\usepackage{microtype}
\usepackage{xspace}
\usepackage{siunitx}
\usepackage{enumitem}
\usepackage{tikz}
\usetikzlibrary{shapes.geometric, arrows.meta, positioning, fit, backgrounds, calc, patterns}
\usepackage{pgfplots}
\pgfplotsset{compat=1.18}
\definecolor{steel_blue}{RGB}{70,130,180}
\definecolor{burnt_orange}{RGB}{204,102,0}
\definecolor{crimson_red}{RGB}{178,34,34}
\usepackage{algorithm}
\usepackage{algpseudocode}
\usepackage{amsthm}
\usepackage[normalem]{ulem} % Review markup: \sout{} for proposed deletions
\usepackage{mathtools}
\usepackage{orcidlink}

\graphicspath{{../figures/}{../results/}}

% ── Theorem-like environments ──
\newtheorem{theorem}{Theorem}[section]
\newtheorem{proposition}[theorem]{Proposition}
\newtheorem{remark}[theorem]{Remark}

% ── Custom commands ──
\newcommand{\modelname}{PC-RadNet}
\newcommand{\ie}{\textit{i.e.}\xspace}
\newcommand{\eg}{\textit{e.g.}\@\xspace}
\newcommand{\etal}{\textit{et al.}\xspace}
\newcommand{\wrt}{w.r.t.\xspace}
\DeclareMathOperator*{\argmin}{argmin}
\DeclareMathOperator*{\argmax}{argmax}

\begin{document}

% ══════════════════════════════════════════════════════════════════
\title{PC-RadNet: A Prompt-Calibrated, Parameter-Efficient Foundation Model for Comprehensive Multi-Task Segmentation, Clinical Assessment and Report Generation in Spine MRI}


\author{%
  \IEEEauthorblockN{Rao Farhat Masood \textsuperscript{\orcidlink{0000-0001-9054-9192}}, \and Imtiaz Ahmad Taj\textsuperscript{\orcidlink{0000-0003-1813-5563}}}%
  \thanks{Rao Farhat Masood and Imtiaz Ahmad Taj are affiliated with the Department of Electrical and Computer Engineering, Capital University of Science and Technology (CUST), Islamabad, Pakistan (e-mail: farhatmasood.fm@gmail.com; imtiaztaj@cust.edu.pk).}\\
  \thanks{Manuscript received - Month -; revised - Month -; accepted - Month -. Date of publication - Month -; date of current version - Month -. (Corresponding author: Rao Farhat Masood.) }
}

\maketitle

% ══════════════════════════════════════════════════════════════════
\begin{abstract}
Lumbar spine MRI analysis requires concurrent anatomical segmentation, disc degeneration grading, disc height quantification, multi-label pathology detection and structured report generation. These tasks are traditionally addressed by separate resource-intensive models that impose spatial annotation requirements at inference. We propose \modelname{}, a lightweight, prompt-calibrated multi-task framework that unifies these capabilities within a single architecture. Built upon a frozen DINOv2 backbone with 2.73 M trainable parameters, \modelname{} introduces a Spine Domain Adapter for radiological feature adaptation, a Prompt Calibrator for multi-conditional feature gating and a multi-task head ensemble for simultaneous prediction of segmentation masks, Pfirrmann grades, disc heights and spinal disorders. A retrieval-augmented template filling (RATF) module generates structured radiology reports from \modelname{} predictions compliant with RSNA guidelines. Trained on multi-planar lumbar \href{https://data.mendeley.com/datasets/p3r4xd2488/1}{Internal Dataset (LSMA-PQR)}, \modelname{} achieves a mean Dice coefficient $\bar{D}=0.861$ and mean IoU $\bar{J}=0.770$ for prompted segmentation, 92.3\% within-one-grade accuracy for Pfirrmann grading, disc height MAE of 0.57 mm and patient-level disorder macro-F1 of 0.57. The RATF report generator attains disorder-level F1 of 0.798, outperforming twelve benchmarked medical LLMs while maintaining a 0.63 GB VRAM footprint at inference. Zero-shot evaluation on \href{https://zenodo.org/records/10159290}{SPIDER} and \href{https://www.kaggle.com/competitions/rsna-2024-lumbar-spine-degenerative-classification/data}{RSNA 2024} confirms cross-institutional transfer of grading performance, with segmentation constrained by acquisition divergence. Code is available at \url{}.
\end{abstract}

\begin{IEEEkeywords}
Lumbar spine MRI, prompt-guided segmentation, multi-task learning, parameter-efficient fine-tuning, DINOv2, LoRA, radiological reporting, Pfirrmann grading
\end{IEEEkeywords}

% ══════════════════════════════════════════════════════════════════
\section{Introduction}
\label{sec:introduction}

Degenerative lumbar spine disease affects up to 40\% of adults at some point in their lifetime \cite{hoy2010epidemiology}, representing a leading cause of disability worldwide. Accurate, timely MRI analysis, including identification of disc herniations, spinal canal stenosis and nerve root compression, is essential for clinical decision-making, yet the process remains labor-intensive and subject to inter-observer variability \cite{lurie2016management}. Radiologist agreement on Pfirrmann grading, the most widely used disc degeneration classification scheme, ranges from $\kappa = 0.56$--$0.76$, corresponding to approximately 70--80\% inter-rater agreement \cite{wilke2006discdegeneration}, establishing a practical ceiling for automation.

Recent deep learning approaches have demonstrated promise in automating individual subtasks: semantic segmentation of vertebral bodies and intervertebral discs \cite{pang2021spineparsenet, chen2021transunet}, Pfirrmann grading of disc degeneration \cite{jamaludin2017spinenet} and structured report generation \cite{jing2019show}. These methods typically require separate architectures for each task, leading to high computational overhead and limited clinical deployability, particularly in resource-constrained clinical environments, including hospital workstations equipped with consumer-grade or entry-level GPU hardware. Multi-task learning \cite{caruana1997multitask} offers a principled solution by sharing representations across related tasks, but existing spine-specific pipelines lack the modular, parameter-efficient design needed for clinical deployment.

Foundation models such as the Segment Anything Model (SAM) \cite{kirillov2023segment} and DINOv2 \cite{oquab2024dinov2} have shown remarkable zero-shot generalization in natural-image domains. More recently, SAM2 \cite{ravi2024sam2} extended prompt-guided segmentation to video and 3D volumes, while MedSAM-2 \cite{medsam2_2024} and SAM2Rad \cite{sam2rad_2024} demonstrated strong performance on medical image segmentation without manual spatial prompts. Dedicated medical segmentation benchmarks, notably nnU-Net \cite{isensee2021nnu}, SwinUNETR \cite{hatamizadeh2022swinunetr} and SAM-Med2D \cite{cheng2023sammed2d}, have established strong performance baselines for anatomical segmentation. Two publicly available lumbar spine benchmarks, SPIDER \cite{van2024spider} and the RSNA 2024 Lumbar Spine Degenerative Conditions challenge \cite{rsna_lumbar_2026}, have recently expanded evaluation resources for this domain. Nevertheless, text-prompted clinical interaction and multi-task spine-specific outputs remain unaddressed by these approaches.

Fine-tuning the full SAM model ($\sim$93 M parameters) is impractical under clinical VRAM constraints. Parameter-efficient fine-tuning (PEFT) strategies, notably Low-Rank Adaptation (LoRA) \cite{mao2024lora_survey}, enable the adaptation of large pre-trained models by training fewer than 5\% of total parameters. Combined with domain-specific modules, this paradigm offers a path to high-performance, resource-efficient medical image analysis.

% \textcolor{red}{\textbf{Novelty Position.}}
Three unresolved gaps in the current literature motivate \modelname{}:
(i) text-prompted medical segmentation methods (\eg SAT-Pro \cite{satpro2025}) rely on natural-image-trained vocabularies that produce near-zero Dice on lumbar MRI; (ii) SAM-based medical methods \cite{ma2024medsam, cheng2023sammed2d, medsam2_2024} require ground-truth bounding boxes or point annotations at inference, which is clinically impractical; and (iii) no prior single architecture simultaneously delivers segmentation, ordinal grading, biometric estimation, disorder detection and structured reporting within one unified forward pass. \modelname{} closes all three gaps through; a learnable anatomy vocabulary with \emph{multi-conditional} feature gating as the first annotation-free text-prompted spine segmentation mechanism (Theorem \ref{thm:expressiveness}, \Cref{sec:prompt_calibrator}); a fully unified five-output clinical pipeline; and a 2.73 M-parameter, 0.63 GB-VRAM deployment footprint suitable for resource-constrained clinical workstations.

% \textbf{Contributions.}
We present \modelname{}, a unified framework that:
\begin{enumerate}
    \item Introduces a \emph{Prompt Calibrator} that learns a closed-vocabulary anatomy embedding with three-way conditional gating (anatomy $\times$ imaging plane $\times$ spinal level), enabling annotation-free text-guided segmentation without bounding-box or point annotations. Theorem \ref{thm:expressiveness} (\Cref{sec:prompt_calibrator}) formally proves its strict expressiveness advantage over all single-conditional variants.
    \item Proposes a zero-identity-initialized Spine Domain Adapter, a rank-32 residual bottleneck with learnable scalar gating that bridges the natural-image to radiological domain gap while preserving pretrained DINOv2 representations at initialization by construction.
    \item Simultaneously produces pixel-level segmentation, ordinal Pfirrmann grade (\Cref{thm:corn}), disc height estimate, nine-class disorder probabilities and structured reports in a single forward pass using only 3.08\% of total parameters and 0.63 GB peak VRAM.
    \item Benchmarks RATF report generation against twelve foundation LLMs on a matched 20-patient protocol, achieving disorder-level F1 $=0.798$ and outperforming all twelve foundation LLMs including Phi-3.5-mini ($0.705$), Med-PaLM ($0.711$), Meditron3-8B ($0.699$), Llama-3.2-3B-Instruct ($0.673$), MedGemma-4B-IT ($0.652$), BioMistral-7B ($0.655$), BioGPT-Large ($0.605$) and LLaVA-Rad-7B ($0.602$), while preserving deterministic RSNA-structured outputs (\Cref{tab:report_comparison}). RATF additionally achieves the highest scores on all NLG metrics including BLEU-1 ($0.221$), METEOR ($0.384$), ROUGE-L ($0.298$) and BERTScore ($0.870$), via a concise clinical prose generator matching radiologist vocabulary (\Cref{tab:nlg_eval}).
    \item Provides a comprehensive evaluation with statistical rigour, including an ablation study (4 variants) with bootstrap 95\% CIs, Wilcoxon signed-rank tests and Bonferroni correction, five direct and seven literature baselines and zero-shot cross-dataset validation on five external datasets (SPIDER, RSNA 2024, SpineWeb, Spine DataSet NEW, scoliosis X-ray).
    \item Contributes the LSMA-PQR dataset comprising 515 patients (4,120 images, seven-class masks, Pfirrmann grades, disc heights, nine-class disorders) and a formal domain adaptation algorithm (\Cref{alg:adaptation}, \Cref{app:generalization}) for extending \modelname{} to seven additional anatomical domains.
\end{enumerate}
The remainder of this paper is organized as follows: \Cref{sec:related} surveys the relevant literature, \Cref{sec:method} details the proposed architecture, \Cref{sec:experiments} describes the experimental setup, \Cref{sec:results} presents quantitative and qualitative results and \Cref{sec:discussion} discusses findings, limitations and future directions.

% ══════════════════════════════════════════════════════════════════
\section{Related Work}
\label{sec:related}

The literature review is organized across five subsections covering anatomical segmentation methods for lumbar spine MRI, automated Pfirrmann grading, parameter-efficient fine-tuning of foundation models, radiological report generation and sentence embeddings for clinical text retrieval.

\subsection{Lumbar Spine Segmentation}
Automated lumbar spine segmentation has evolved from atlas-based methods \cite{zukic2014robust} through U-Net variants \cite{ronneberger2015unet} to transformer-based architectures.
SpineParseNet \cite{pang2021spineparsenet} proposed a two-stage coarse-to-fine approach for multi-class segmentation, while TransUNet \cite{chen2021transunet} combined CNN features with transformer encoders.
nnU-Net \cite{isensee2021nnu}, a self-configuring segmentation framework, has become the de facto benchmark for medical image segmentation across modalities and SwinUNETR \cite{hatamizadeh2022swinunetr} demonstrated that Swin Transformers can match or exceed CNN-based methods on volumetric tasks.

Recent SAM-based methods \cite{ma2024medsam} adapt the Segment Anything Model to medical imaging via prompt engineering or lightweight adapters, but typically require bounding-box prompts derived from ground-truth annotations, an impractical assumption in clinical deployment.
SAM-Med2D \cite{cheng2023sammed2d} fine-tunes SAM on 4.6M medical images with point and box prompts.
UniverSeg \cite{butoi2023universeg} and EfficientSAM \cite{xiong2024efficientsam} explore few-shot and efficient alternatives.
MedSAM-2 \cite{medsam2_2024} adapted SAM2 for universal 2D/3D medical image segmentation across CT, MRI and PET modalities, establishing new performance baselines.
SAM2Rad \cite{sam2rad_2024} demonstrated automatic prompt generation for SAM/SAM2 on medical images using as few as 10 labeled examples, directly related to our no-oracle-prompt objective.
MedUniSeg \cite{meduniseg_2024} proposed prompt-driven universal segmentation, while DAPSAM \cite{dapsam_2025} introduced domain-adaptive prototypes for generalizable medical image segmentation.
SAT-Pro \cite{satpro2025} recently demonstrated large-vocabulary (497-category) text-prompted 3D medical image segmentation, outperforming nnU-Net on several anatomical regions; however, SAT-Pro targets general-purpose 3D segmentation without the multi-task clinical outputs (grading, height, disorders, reporting) that define \modelname{}'s scope.
ATM-Net \cite{atmnet2025} proposes anatomy-aware text-guided multi-modal fusion specifically for lumbar spine segmentation, complementing our approach with an alternative text-to-feature fusion strategy.
Notably, none of these approaches addresses the simultaneous multi-task objectives of segmentation, ordinal grading, disc height regression and structured report generation from a single unified pass, as proposed in \modelname{}.

\subsection{Disc Degeneration Grading}
Pfirrmann grading \cite{pfirrmann2001magnetic} provides an ordinal scale (I--V) for disc degeneration severity. As shown in \Cref{figure1_pfirrmann}, Grade I represents a healthy disc with homogeneous bright signal and clear distinction between nucleus and annulus, while Grade V indicates severe degeneration with collapsed disc space and loss of signal. SpineNet \cite{jamaludin2017spinenet} pioneered automated grading with a multi-task CNN, while subsequent work incorporated ordinal regression losses \cite{cao2020rank} and uncertainty estimation \cite{liawrungrueang2023discdegeneration}. Our approach treats grading as cumulative ordinal regression with the CORN loss \cite{shi2022corn}, enabling principled uncertainty quantification through threshold calibration. A persistent challenge in automated Pfirrmann grading is severe class imbalance: Grades II--III dominate clinical populations while Grades I and V are rare \cite{jamaludin2017spinenet}. Recent meta-analyses report inter-observer agreement of $\kappa = 0.56$--$0.76$ for Pfirrmann grading \cite{wilke2006discdegeneration}, indicating that even human experts frequently disagree on adjacent grades, a ceiling that should contextualize automated performance.

\begin{figure}[!htbp]
    \includegraphics[width=\columnwidth]{figure1.pdf}
    \caption{Pfirrmann Grading System for Lumbar Disc Degeneration.}
    \label{figure1_pfirrmann}
\end{figure}

\subsection{Parameter-Efficient Foundation Model Adaptation}
LoRA \cite{mao2024lora_survey} injects low-rank update matrices into frozen transformer weights, achieving comparable performance to full fine-tuning at a fraction of the parameter cost. Medical imaging applications have adopted LoRA for radiology report generation \cite{chen2024chexagent} and pathology analysis \cite{lu2024visual}. Med-SA \cite{wu2025medsamadapter} adapts SAM for medical images using adapter modules, while SAMed \cite{zhang2023customized} fine-tunes SAM's image encoder with LoRA. Our work extends this paradigm by coupling LoRA with a novel multi-conditional prompt calibration mechanism (anatomy $\times$ plane $\times$ spinal level) and multi-task heads, achieving unified segmentation-grading-reporting in a single architecture. We note that class-conditional feature modulation has been explored since conditional BatchNorm \cite{de2017modulating}; the novelty of our Prompt Calibrator lies not in the vocabulary lookup per se but in the \emph{multi-conditional} calibration (three-way anatomy--plane--level interaction) combined with the multi-task PEFT efficiency story. Subsequent PEFT advances include AdaLoRA \cite{zhang2023adalora}, which dynamically allocates rank budget across layers via singular value decomposition and QLoRA \cite{dettmers2023qlora}, which combines 4-bit quantization with LoRA to further reduce memory footprint. DoRA \cite{liu2024dora} decomposes weight updates into magnitude and direction components, consistently outperforming vanilla LoRA by 2--5\% across language and vision tasks. \modelname{} adopts vanilla LoRA ($r=8$) for simplicity and reproducibility; we identify DoRA and AdaLoRA as promising upgrade paths in \Cref{sec:discussion}.

\subsection{Radiological Report Generation}
Automated report generation spans template-based \cite{li2018hybrid}, retrieval-augmented \cite{chen2023ptunifier} and end-to-end neural approaches \cite{chen2020generating}. For structured spine reporting, template-based methods offer deterministic, clinically precise outputs. We adopt a hybrid approach: template-based generation augmented with FAISS \cite{johnson2019billion} nearest-neighbour retrieval to select contextually relevant exemplar reports. Recent large multimodal models, including MAIRA-2 \cite{bannur2024maira2}, which generates grounded radiology reports with pixel-level visual attribution, RadFM \cite{wu2025radfm}, a generalist 2D/3D radiology foundation model, LLaVA-Med \cite{li2023llava_med}, a biomedical large vision-language assistant and MedGemma \cite{medgemma2025}, a medical-domain multimodal model built on Gemma 3 with SigLIP vision encoding (4.97B parameters), demonstrate the potential of end-to-end neural report generation. Med-PaLM 2 \cite{singhal2023medpalm2} further illustrates the ceiling of large-scale medical LLMs, achieving expert-level performance on medical question answering benchmarks; however, no public pre-trained weights are available and the open-source re-implementation initializes with random parameters, limiting reproducibility. However, these systems require substantial compute (billions of parameters) and lack structured, template-compliant outputs aligned with RSNA reporting guidelines \cite{rsna2024templates}. Concurrently, the SPINE framework \cite{spine_framework_2026} introduced a segmentation-guided pipeline for spinal MRI report generation; \modelname{} differs by integrating prompting, multi-task heads and retrieval-augmented template filling within a parameter-efficient single-model architecture. \modelname{} targets the complementary regime: deterministic, clinically structured reports at 2.73 M trainable parameters.

\subsection{Sentence Embeddings for Clinical Text Retrieval}
General-purpose sentence encoders such as all-MiniLM-L6-v2 \cite{wang2020minilm} achieve strong performance on open-domain retrieval but may underperform on clinical text due to domain vocabulary mismatch \cite{jin2023medcpt}. MedCPT \cite{jin2023medcpt} applies contrastive pre-training on 255,000 clinician-validated query--article pairs from PubMed, producing embeddings specifically calibrated for biomedical retrieval. ClinicalBERT \cite{alsentzer2019clinicalbert}, trained on MIMIC-III discharge summaries, better captures the syntactic conventions of radiological reports. \modelname{} currently uses all-MiniLM-L6-v2 for the FAISS retrieval index; replacing this with MedCPT is identified as a priority upgrade in \Cref{sec:discussion}. Based on the identified gaps, the following section details how \modelname{} attempts to address the challenges through its modular architecture.


\begin{figure*}[!t]
	\centering
	\includegraphics[width=\textwidth]{figure2.pdf}
	\caption{%
Architecture of \modelname{}. A frozen DINOv2 ViT-B/14 backbone extracts spatial tokens from the input MRI, adapted via LoRA (teal) across all transformer blocks. A Spine Domain Adapter (purple) bridges the natural-to-medical domain gap using a learnable residual bottleneck. Concurrently, a Prompt Calibrator encodes anatomical, view ($\pi$), and spinal level ($\lambda$) embeddings into a calibration vector ($\mathbf{c}$) that modulates the upsampled features in the Segmentation Decoder via sigmoid gating. The network bifurcates into dense and global prediction pathways: the decoder resolves prompted and multi-class segmentation masks (axial and sagittal), while mean-pooled global features ($\overline{\mathbf{f}}$) are processed by specific task heads (blue) to predict ordinal Pfirrmann grades ($\hat{g}$), disc heights ($\hat{h}$), and multi-label disorders ($\hat{\mathbf{d}}$). These predictions culminate in the generation of structured radiology reports via an integrated RATF module (orange).}
	\label{fig:architecture}
\end{figure*}

% ══════════════════════════════════════════════════════════════════
\section{Method}
\label{sec:method}

We present the six components of \modelname{} encompassing the frozen DINOv2 backbone with LoRA adaptation, the Spine Domain Adapter, the Prompt Calibrator, the multi-task prediction heads, the composite loss function and the report generator. \Cref{fig:architecture} illustrates the overall architecture of our proposed \modelname{} framework. The system is designed to process a multi-planar lumbar MRI image, $\mathbf{x} \in \mathbb{R}^{3 \times H \times W}$, alongside a clinical text prompt, $p$, selected from a predefined vocabulary of 20 anatomical terms (\eg, ``intervertebral disc''). By integrating these multimodal inputs, \modelname{} jointly predicts a comprehensive suite of clinical markers, including prompt-guided binary segmentation masks ($\hat{m}_p$) and multi-class segmentation logits for regional anatomy across axial and sagittal views. Furthermore, the model outputs an ordinal Pfirrmann degeneration grade ($\hat{g} \in \{1,\ldots,5\}$), a continuous disc height estimate ($\hat{h} \in \mathbb{R}_+$) and multi-label classifications for spinal disorders. These unified outputs ultimately feed into a RATF module to synthesize standardized, RSNA-compliant radiological reports. The subsequent subsections detail the core components of this architecture, including the multimodal representation learning, prompt calibration and the multi-task decoding pathways.

% \input{fig1_architecture_tikz}


\subsection{Frozen Backbone with LoRA}
We adopt DINOv2 ViT-Base/14 \cite{oquab2024dinov2}, pre-trained via self-supervision on 142 M natural images, as the frozen vision encoder. With a patch size of 14 and input resolution $384 \times 384$, the backbone produces $\lfloor 384/14 \rfloor^2 = 27^2 = 729$ spatial tokens of dimension $d=768$, handled internally by the \texttt{timm} patch embedding convolution. Since $384 \bmod 14 = 6$, the patch grid implicitly discards a 6-pixel boundary strip; at the typical lumbar MRI pixel spacing of $\sim$0.5 mm, this corresponds to $\sim$3 mm of peripheral context, which is negligible relative to the central anatomy of interest. We inject LoRA \cite{mao2024lora_survey} adapters into the fused QKV projection, output projection and both FFN layers (\texttt{fc1}, \texttt{fc2}) of all 12 transformer blocks:

\begin{equation}
    \mathbf{W}' = \mathbf{W}_0 + \frac{\alpha}{r} \mathbf{B}\mathbf{A},
    \quad \mathbf{A} \in \mathbb{R}^{r \times d_{\text{in}}}, \; \mathbf{B} \in \mathbb{R}^{d_{\text{out}} \times r},
    \label{eq:lora}
\end{equation}
where $r=8$, $\alpha=16$ and dropout $0.05$.
Applying LoRA to QKV ($768 \to 2304$), output projection ($768 \to 768$) and FFN ($768 \to 3072 \to 768$) across 12 blocks introduces 1.18 M trainable parameters (1.37\% of the 86 M backbone).

\subsection{Spine Domain Adapter}
\label{sec:spine_adapter}
The Spine Domain Adapter bridges the domain gap between natural-image and medical features via a residual bottleneck with learnable scaling and zero-initialized up-projection. Let $\mathbf{f} \in \mathbb{R}^{N \times d}$ denote the token features extracted by the DINOv2 backbone ($N=729$, $d=768$).
The adapter applies a residual bottleneck projection:
\begin{equation}
    \mathbf{f}_{\text{spine}} = \mathbf{f} + \alpha_s \cdot \mathbf{W}_{\text{up}}   \phi\!\left(\mathbf{W}_{\text{down}} \cdot \text{LN}\!\left(\mathbf{f}\right)\right),
    \label{eq:sda}
\end{equation}
where $\mathbf{W}_{\text{down}} \!\in\! \mathbb{R}^{d_r \times d}$, $\mathbf{W}_{\text{up}} \!\in\! \mathbb{R}^{d \times d_r}$ with adapter rank $d_r=32$, $\phi(\cdot) = \text{GELU}(\cdot)$, $\text{LN}(\cdot)$ is Layer Normalization and $\alpha_s$ is a learnable scaling parameter initialized to 0.1.
% Three design choices are anatomically motivated:
% \begin{enumerate}[nosep, leftmargin=*]
%     \item \textbf{Zero-initialized up-projection} ($\mathbf{W}_{\text{up}} = \mathbf{0}$ at $t=0$): ensures the adapter behaves as an identity mapping at initialization, \textcolor{red}{preserving pre-trained representations and preventing catastrophic forgetting}.
%     \item \textbf{Residual connection with learnable scale}: \textcolor{red}{the scalar $\alpha_s$ controls the magnitude of domain adaptation at each training step, accommodating inter-scanner variability (1.5T vs. 3T, varying coil configurations) without overriding pre-trained anatomical priors.}
%     \item \textbf{Low-rank bottleneck} ($d_r=32 \ll d=768$): \textcolor{red}{restricts adaptation to $d_r$ principal directions out of $d$, reducing overfitting on the training set ($n=2{,}880$) while preserving the remaining feature structure.}
% \end{enumerate}

The proposed methodology incorporates three anatomically motivated design choices to optimize domain adaptation while retaining robust foundational features. First, a zero-initialized up-projection ($\mathbf{W}_{\text{up}} = \mathbf{0}$ at $t=0$) guarantees that the adapter initially functions as an identity mapping, thereby preserving pre-trained representations and preventing representation drift. Second, a residual connection modulated by a learnable scale parameter ($\alpha_s$) controls the magnitude of adaptation at each training step, accommodating inter-scanner variability such as differences between 1.5T and 3T magnetic field strengths or diverse coil configurations, without overriding established pre-trained anatomical priors. Finally, the implementation of a low-rank bottleneck ($d_r=32 \ll d=768$) restricts the adaptation process to a constrained number of principal directions. This targeted dimensionality reduction significantly mitigates the risk of overfitting on the training set ($n=2,880$) while preserving the integrity of the broader feature structure.

The adapted features retain the spatial token structure ($\mathbf{f}_{\text{spine}} \in \mathbb{R}^{N \times d}$), enabling direct integration with the downstream Prompt Calibrator and Segmentation Decoder. The adapter introduces $2 \cdot d \cdot d_r + (d_r + d) + 2d + 1 = 51,489$ additional parameters, including Layer Normalization weights and biases.

\subsection{Prompt Calibrator}
\label{sec:prompt_calibrator}
The Prompt Calibrator conditions the segmentation decoder on clinical text prompts from a vocabulary of 20 anatomical terms, eliminating the need for bounding-box or point annotations required by prior methods \cite{kirillov2023segment, cheng2023sammed2d, ma2024medsam}.

\subsubsection{Prompt Encoding}
Given a clinical text prompt $p$ (\eg ``intervertebral disc'', ``spinal canal''), the Prompt Calibrator encodes it through a \emph{learnable anatomy vocabulary} rather than a frozen text encoder, thereby avoiding domain mismatch of general-purpose encoders on medical terminology:
\begin{equation}
    \mathbf{e}_p = \text{Emb}_{\text{anat}}(\text{idx}(p)), \quad \mathbf{e}_p \in \mathbb{R}^{d},
    \label{eq:anat_embed}
\end{equation}
where $\text{Emb}_{\text{anat}}$ is a learnable embedding matrix with 20 entries (one per anatomy term in the prompt bank) and $\text{idx}(\cdot)$ maps prompt text to the corresponding vocabulary index. A two-layer MLP projects the anatomy embedding to the calibration space:
\begin{equation}
    \mathbf{t}_p = \mathbf{W}_2^{\text{proj}}   \phi\!\left(\mathbf{W}_1^{\text{proj}}   \mathbf{e}_p + \mathbf{b}_1^{\text{proj}}\right) + \mathbf{b}_2^{\text{proj}}, \quad \mathbf{t}_p \in \mathbb{R}^{d_c},
    \label{eq:text_proj}
\end{equation}
with calibration dimension $d_c=256$.

\subsubsection{Multi-Conditional Calibration}
The calibration vector is derived by concatenating three conditioning signals, namely anatomy, imaging plane and spinal level, followed by a learned non-linear fusion:
\begin{equation}
    \mathbf{c} = \text{MLP}\!\left(\left[\mathbf{t}_p \;\|\; \mathbf{v}_\pi \;\|\; \boldsymbol\ell_\lambda\right]\right), \quad \mathbf{c} \in \mathbb{R}^{d_c},
    \label{eq:calibrator}
\end{equation}
where $\mathbf{v}_\pi = \text{Emb}_{\text{view}}(\pi)$ encodes the imaging plane ($\pi \in \{0, 1\}$ for axial, sagittal; $\mathbf{v}_\pi \in \mathbb{R}^{d_c}$), $\boldsymbol\ell_\lambda = \text{Emb}_{\text{level}}(\lambda)$ encodes the disc level ($\lambda \in \{0, \ldots, 5\}$; $\boldsymbol\ell_\lambda \in \mathbb{R}^{d_c}$) and $[\cdot \| \cdot]$ denotes concatenation. Both $\text{Emb}_{\text{view}}$ and $\text{Emb}_{\text{level}}$ produce $d_c=256$-dimensional embeddings, yielding an input dimensionality of $3d_c=768$ for the fusion MLP. The MLP consists of two linear layers with GELU activation and dropout (rate $0.1$): $\mathbb{R}^{3d_c} \to \mathbb{R}^{d_c} \to \mathbb{R}^{d_c}$.

This design integrates the \emph{anatomy prompt} $\mathbf{t}_p$, identifying the target structure; the \emph{plane embedding} $\mathbf{v}_\pi$, distinguishing axial morphology (disc cross-section, posterior elements) from sagittal morphology (vertebral column, lordotic curvature); and the \emph{level embedding} $\boldsymbol\ell_\lambda$, encoding level-specific variations (\eg L5/S1 transitional anatomy); hence, encoding three clinically significant anatomical priors. Since $\lambda \in \{0,\ldots,5\}$ encodes 3 disc levels $\times$ 2 planes, partial redundancy with $\mathbf{v}_\pi$ exists; the concatenated input allows the fusion MLP to learn joint level$\times$plane interactions.

\subsubsection{Multiplicative Feature Gating}
\label{sec:gating}
The adapted spatial tokens are projected to the calibration dimension: $\mathbf{u}_i = \mathbf{W}_{\text{tok}} \mathbf{f}_{\text{spine},i} + \mathbf{b}_{\text{tok}}$, where $\mathbf{W}_{\text{tok}} \in \mathbb{R}^{d_c \times d}$, yielding $\mathbf{u}_i \in \mathbb{R}^{d_c}$. Here, we used $\mathbf{u}$ to distinguish projected tokens from the input image $\mathbf{x}$.
The calibration vector $\mathbf{c}$ modulates the projected features via a sigmoid-gated multiplicative mechanism:
\begin{equation}
    \hat{\mathbf{u}}_i = \mathbf{u}_i \odot \left(\mathbf{1} + \sigma\!\left(\mathbf{W}_g   \mathbf{c}\right)\right), \quad i \in \{1, \ldots, N\},
    \label{eq:gate}
\end{equation}
where $\mathbf{u}_i \in \mathbb{R}^{d_c}$ is the $i$-th projected spatial token, $\sigma(\cdot)$ is the element-wise sigmoid function and $\mathbf{W}_g \in \mathbb{R}^{d_c \times d_c}$. The additive identity term yields a baseline scaling of $1.5\times$ when $\mathbf{c}=\mathbf{0}$ (since $\sigma(\mathbf{0})=0.5$), bounding the per-channel range to $(1.0, 2.0)$: channels relevant to the queried anatomy are amplified up to $2\times$, while irrelevant channels receive near-unit scaling. This FiLM-style (Feature-wise Linear Modulation) \cite{perez2018film} multiplicative gating acts as soft channel-wise attention, where the prompt selects which feature dimensions to emphasize uniformly across all spatial locations. The formulation does not attenuate channels below their input magnitude; bidirectional modulation via $2\sigma(\mathbf{W}_g \mathbf{c})$ is left for future work.

% \textcolor{red}{\textbf{Interpretability and reproducibility.}
The gate values $\sigma(\mathbf{W}_g \mathbf{c})$ provide per-channel attention scores for a given prompt, enabling post-hoc interpretability without external attribution methods.
The mapping is fully deterministic ensuring identical inputs ($\mathbf{x}$, $p$, $\pi$, $\lambda$) always produce the same segmentation mask, ensuring clinical reproducibility \cite{rsna2024templates}.

To formally justify this architectural design, we demonstrate that conditioning the feature gate on the joint concatenation of these three embeddings provides strictly greater representational capacity than any isolated conditional formulation.

\begin{theorem}
\label{thm:expressiveness}
Let $\mathcal{G}_{\mathcal{A}}$, $\mathcal{G}_{\mathcal{P}}$, and $\mathcal{G}_{\mathcal{L}}$ denote the hypothesis classes of channel-wise gate functions conditioned exclusively on anatomy, imaging plane, and disc level, respectively. Let $\mathcal{G}_{\mathcal{APC}}$ denote the multi-conditional hypothesis class defined by the Prompt Calibrator (Eq. \ref{eq:calibrator}). Then, $\mathcal{G}_{\mathcal{A}} \cup \mathcal{G}_{\mathcal{P}} \cup \mathcal{G}_{\mathcal{L}} \subsetneq \mathcal{G}_{\mathcal{APC}}$.
\end{theorem}

\textbf{Proof.}
Let the input embeddings be denoted as $\mathbf{t}_p \in \mathbb{R}^{d_t}$, $\mathbf{v}_\pi \in \mathbb{R}^{d_v}$, and $\boldsymbol\ell_\lambda \in \mathbb{R}^{d_l}$. The multi-conditional gate $g \in \mathcal{G}_{\mathcal{APC}}$ is parameterized by a Multi-Layer Perceptron (MLP) operating on the concatenated joint space $\mathbf{x} = [\mathbf{t}_p \mathbin\Vert \mathbf{v}_\pi \mathbin\Vert \boldsymbol\ell_\lambda]$.

(\textit{Containment} $\supseteq$): Any isolated gating function $g_A \in \mathcal{G}_{\mathcal{A}}$ depends exclusively on $\mathbf{t}_p$. This function can be exactly recovered within $\mathcal{G}_{\mathcal{APC}}$ by setting the weight matrices corresponding to the orthogonal subspaces of $\mathbf{v}_\pi$ and $\boldsymbol\ell_\lambda$ in the first linear layer of the fusion MLP to zero. By symmetry, this holds for $\mathcal{G}_{\mathcal{P}}$ and $\mathcal{G}_{\mathcal{L}}$, thereby establishing $\mathcal{G}_{\mathcal{A}} \cup \mathcal{G}_{\mathcal{P}} \cup \mathcal{G}_{\mathcal{L}} \subseteq \mathcal{G}_{\mathcal{APC}}$.

(\textit{Strict Containment} $\subsetneq$): Consider a target gating function $g^*$ that requires non-separable cross-conditional interactions, such as $g^*(\mathbf{t}_p, \mathbf{v}_\pi) \propto \mathbf{t}_p^\top \mathbf{W} \mathbf{v}_\pi$. No function residing in the union of the single-conditional classes can represent this interaction, as elements of $\mathcal{G}_{\mathcal{A}} \cup \mathcal{G}_{\mathcal{P}} \cup \mathcal{G}_{\mathcal{L}}$ map strictly from isolated embedding spaces. Conversely, the fusion MLP in $\mathcal{G}_{\mathcal{APC}}$ projects from the concatenated joint space, inherently providing the capacity to model these high-order dependencies. Therefore, the containment is strict.

This theoretical capacity explains the substantial segmentation improvements observed empirically when the network is resolving cross-plane morphological ambiguities.


\subsection{Multi-Task Heads}
\label{sec:heads}

The gated features $\hat{\mathbf{U}} \in \mathbb{R}^{N \times d_c}$ are processed through a shared convolutional decoder and four task-specific heads.

\subsubsection{Segmentation Decoder}
The spatial tokens $\hat{\mathbf{U}}$ are reshaped to a 2D feature map $\mathbf{Z}_0 \in \mathbb{R}^{d_c \times h \times w}$ where $h=w=27$, then progressively upsampled through four transposed convolution stages:
\begin{equation}
    \mathbf{Z}_{l+1} = \phi\!\left(\text{BN}\!\left(\text{ConvT}_{k=2, s=2}^{(l)}(\mathbf{Z}_l)\right)\right), \quad l \in \{0,1,2,3\},
    \label{eq:decoder}
\end{equation}
with channel dimensions $256 \to 128 \to 64 \to 32 \to 16$ and spatial dimensions $27 \to 54 \to 108 \to 216 \to 432$, followed by bilinear interpolation to $384 \times 384$.
Three separate $1 \times 1$ convolution heads produce plane-specific and prompted segmentation logits:

% \begin{align}
%     \hat{\mathbf{S}}_{\text{ax}} &= \text{Conv}_{1\times 1}^{16 \to 4}(\mathbf{Z}_4) \in \mathbb{R}^{4 \times 384 \times 384}, \label{eq:axial_head}\\
%     \hat{\mathbf{S}}_{\text{sag}} &= \text{Conv}_{1\times 1}^{16 \to 5}(\mathbf{Z}_4) \in \mathbb{R}^{5 \times 384 \times 384}, \label{eq:sag_head}\\
%     \hat{m}_p &= \text{Conv}_{1\times 1}^{16 \to 1}(\mathbf{Z}_4) \in \mathbb{R}^{1 \times 384 \times 384}. \label{eq:prompt_head}
% \end{align}

\begin{equation} \label{eq:segmentation_heads}
\begin{aligned}
    \hat{\mathbf{S}}_{\text{ax}} &= \text{Conv}_{1\times 1}^{16 \to 4}(\mathbf{Z}_4) \in \mathbb{R}^{4 \times 384 \times 384}, \\
    \hat{\mathbf{S}}_{\text{sag}} &= \text{Conv}_{1\times 1}^{16 \to 5}(\mathbf{Z}_4) \in \mathbb{R}^{5 \times 384 \times 384}, \\
    \hat{m}_p &= \text{Conv}_{1\times 1}^{16 \to 1}(\mathbf{Z}_4) \in \mathbb{R}^{1 \times 384 \times 384}.
\end{aligned}
\end{equation}
Here, $\mathbf{Z}_4$ represents the final decoded feature map with 16 channels, and $\text{Conv}_{1\times 1}^{a \to b}$ denotes a $1 \times 1$ convolutional layer mapping $a$ input channels to $b$ output channels. The variables $\hat{\mathbf{S}}_{\text{ax}}$ and $\hat{\mathbf{S}}_{\text{sag}}$ represent the predicted multi-class segmentation logits for the axial and sagittal views respectively, while $\hat{m}_p$ represents the targeted binary segmentation mask generated in response to the clinical prompt. The plane selector $\pi$ determines which multi-class head is active during training and inference: $\hat{\mathbf{S}} = \hat{\mathbf{S}}_{\text{ax}}$ if $\pi=0$, else $\hat{\mathbf{S}}_{\text{sag}}$. The prompted head $\hat{m}_p$ produces binary logits for the queried anatomy regardless of the active plane.

\subsubsection{Grade Head (CORN Ordinal Regression)}
Pfirrmann grading is an inherently ordinal task: Grade III is ``between'' Grade II and Grade IV in a meaningful clinical sense. Standard classification losses ignore this ordering, potentially producing Grade I predictions for Grade V images. We adopt CORN (Conditional Ordinal Regression for Neural networks) \cite{shi2022corn}, which factorizes the ordinal distribution into $K{-}1=4$ conditional binary probabilities:
\begin{equation}
    P(Y > k \mid Y > k{-}1) = \sigma\!\left(\theta_k(\overline{\mathbf{f}})\right), \quad k \in \{1,2,3,4\}
    \label{eq:corn}
\end{equation}
where the conditioning $Y > k{-}1$ is equivalent to $Y \geq k$ for integer-valued grades, $\overline{\mathbf{f}} = \frac{1}{N}\sum_{i=1}^{N} \mathbf{f}_{\text{spine},i} \in \mathbb{R}^d$ is the mean-pooled domain-adapted feature and $\theta_k$ is a shared-backbone MLP ($768 \to 256 \to 64$) with task-specific output neurons.
The predicted grade is:
\begin{equation}
    \hat{g} = 1 + \sum_{k=1}^{K-1} \mathbf{1}\!\left[\prod_{j=1}^{k} \sigma(\theta_j) > 0.5\right].
    \label{eq:grade_pred}
\end{equation}
The CORN loss guarantees monotonic cumulative probabilities $P(Y \ge k) \ge P(Y \ge k{+}1)$ without explicit constraints and the conditional formulation decorrelates adjacent thresholds, reducing gradient interference during training.
This is particularly important for the lumbar spine grading task where extreme class imbalance (Grade I: $n=36$; Grade V: $n=60$ in the training set as compared to Grade II: $n=1164$) causes standard cross-entropy to collapse to majority-class predictions.

\subsubsection{Height and Disorder Heads}
Disc height is estimated via regression on the mean-pooled features:
\begin{equation}
    \hat{h} = \text{ReLU}\!\left(\mathbf{w}_h^\top   \phi\!\left(\mathbf{W}_h   \overline{\mathbf{f}} + \mathbf{b}_h\right)\right) \in \mathbb{R}_+, \quad \text{(mm)}
    \label{eq:height}
\end{equation}
using a two-layer MLP ($768 \to 128 \to 1$) with ReLU to ensure non-negativity. Global mean pooling aggregates features from all structures; however, disc-specific supervision from grading and disorder labels implicitly biases the pooled representation towards disc-level information. Region-of-interest pooling is discussed as a refinement in \Cref{sec:discussion}.
Multi-label disorder classification for 9 binary indicators $\hat{\mathbf{d}} = \sigma\!\left(\mathbf{W}_d   \phi(\mathbf{W}'_d   \overline{\mathbf{f}} + \mathbf{b}'_d) + \mathbf{b}_d\right) \in [0,1]^9$ is performed by a two-layer MLP ($768 \to 256 \to 9$) with sigmoid outputs.

\subsection{Loss Function}
\label{sec:loss}

To optimize the diverse clinical tasks jointly, \modelname{} is trained end-to-end using a composite multi-task objective. The total training loss is defined as a weighted linear combination of six task-specific components:
\begin{equation}
    \mathcal{L} = \lambda_\text{s}\mathcal{L}_\text{seg} + \lambda_\text{p}\mathcal{L}_\text{prompt} + \lambda_\text{b}\mathcal{L}_\text{bound} + \lambda_\text{g}\mathcal{L}_\text{grade} + \lambda_\text{h}\mathcal{L}_\text{ht} + \lambda_\text{d}\mathcal{L}_\text{dis},
    \label{eq:loss}
\end{equation}
where the empirical weights are set to $(\lambda_\text{s}, \lambda_\text{p}, \lambda_\text{b}, \lambda_\text{g}, \lambda_\text{h}, \lambda_\text{d}) = (1.0, 1.0, 0.2, 0.3, 0.2, 0.3)$. This weighting strategy prioritizes the dense, pixel-level segmentation tasks while leveraging the global clinical predictions as auxiliary regularizers.

The regional multi-class segmentation loss ($\mathcal{L}_\text{seg}$) and the binary prompted mask loss ($\mathcal{L}_\text{prompt}$) are both formulated as a combination of class-balanced binary cross-entropy (BCE) and a soft Dice loss. For each foreground class $c \in \{1, \ldots, C\}$, the segmentation loss is computed as:
\begin{equation}
\begin{split}
    \mathcal{L}_\text{seg} = \frac{1}{C}\sum_{c=1}^{C} \Bigl[
    &-\frac{1}{HW}\!\sum_{j}\bigl[w_c  y_{c,j} \log \hat{p}_{c,j} \\
    &\quad+ (1{-}y_{c,j}) \log (1{-}\hat{p}_{c,j})\bigr] \\
    &+ 1 - \frac{2\sum_j y_{c,j}\hat{p}_{c,j} + \epsilon}{\sum_j y_{c,j} + \sum_j \hat{p}_{c,j} + \epsilon}\Bigr].
\end{split}
    \label{eq:seg_loss}
\end{equation}
To mitigate severe anatomical class imbalance—such as the axial spinal canal representing less than 0.25\% of the total pixels—we apply inverse-frequency weights $w_c = N_\text{total} / (C \cdot N_c)$ as a \texttt{pos\_weight} to the BCE term (\eg, $w_\text{SC} \approx 90.0$). The prompted loss, $\mathcal{L}_\text{prompt}$, follows an identical formulation applied to the prompted binary head $\hat{m}_p$.

While often treated empirically as a heuristic penalty, we formalize this inverse-frequency weighting as a structural dataset calibration.

\begin{theorem}
\label{thm:weighting}
Let $P_{\mathcal{D}}(c) = N_c/N_\text{total}$ denote the empirical data distribution for class $c$. Minimizing the weighted empirical loss $\mathcal{L}_w = \sum_{c=1}^{C} w_c \mathcal{L}_{\mathrm{BCE},c}$ with inverse-frequency weights $w_c = 1/(C \cdot P_{\mathcal{D}}(c))$ is strictly equivalent to minimizing the expected unweighted loss under a uniform target distribution $P_U(c) = 1/C$.
\end{theorem}

\textbf{Proof.}
By the definition of the expected loss under the target uniform distribution $P_U$:
\begin{align}
    \mathbb{E}_{P_U}\bigl[\mathcal{L}_{\mathrm{BCE}}\bigr] 
    &= \sum_{c=1}^{C} P_U(c) \mathcal{L}_{\mathrm{BCE},c} \nonumber \\
    &= \sum_{c=1}^{C} P_{\mathcal{D}}(c) \left( \frac{P_U(c)}{P_{\mathcal{D}}(c)} \right) \mathcal{L}_{\mathrm{BCE},c}.
\end{align}
Substituting $P_U(c) = 1/C$, the likelihood ratio becomes $P_U(c)/P_{\mathcal{D}}(c) = 1/(C \cdot P_{\mathcal{D}}(c))$, which exactly matches our definition of $w_c$. Therefore:
\begin{equation}
    \mathbb{E}_{P_U}\bigl[\mathcal{L}_{\mathrm{BCE}}\bigr] = \sum_{c=1}^{C} P_{\mathcal{D}}(c) w_c \mathcal{L}_{\mathrm{BCE},c} = \mathbb{E}_{P_{\mathcal{D}}}\bigl[w_c \mathcal{L}_{\mathrm{BCE}}\bigr].
\end{equation}

This confirms that $w_c$ acts as the exact importance sampling weight required to correct the domain shift from the imbalanced empirical prior to a balanced prior, fundamentally enabling the network to recover extreme minority classes like the axial spinal canal without collapsing.

To enforce structural fidelity, a boundary emphasis loss penalizes predictions that misalign with anatomical edges. This is computed using a Sobel filter applied to the predicted mask and ground truth: $\mathcal{L}_\text{bound} = \|\nabla_{\text{Sobel}} \hat{m}_p - \nabla_{\text{Sobel}} m_p\|_1$. The structural impact of amplifying this penalty (\eg, increasing $\lambda_\text{b}$ to 1.0) is further explored in Ablation A3 (\Cref{sec:ablation}).

For Pfirrmann degeneration grading, we treat the problem as an ordinal regression task to preserve the inherent hierarchy of severity levels. Utilizing the Conditional Ordinal Regression Neural Network (CORN) framework \cite{shi2022corn}, the grading loss ($\mathcal{L}_\text{grade}$) is decomposed into $K{-}1$ conditional binary classification tasks:
\begin{equation}
\begin{split}
    \mathcal{L}_\text{grade} = -\sum_{k=1}^{K-1} \frac{1}{|\mathcal{S}_k|}\sum_{i \in \mathcal{S}_k} \bigl[&y_{ik}^* \log \sigma(\theta_k) \\
    &+ (1 {-} y_{ik}^*) \log (1 {-} \sigma(\theta_k))\bigr],
\end{split}
    \label{eq:corn_loss}
\end{equation}
where $\mathcal{S}_k = \{i : g_i \geq k\}$ denotes the conditioning set of samples that have at least rank $k$, $y_{ik}^* = \mathbf{1}[g_i > k]$ is the binary indicator and $\theta_k$ represents the $k$-th ordinal logit.


The grade head design rests on the structural guarantee that CORN (\cref{thm:corn}) produces monotonically consistent cumulative grade probabilities without any explicit constraint enforcement, as proved below.

\begin{theorem}
\label{thm:corn}
For any finite parameter configuration $\theta$, the cumulative probabilities $\hat{Q}_k \triangleq P_\theta(Y \geq k) = \prod_{j=1}^{k-1}\sigma\!\left(\theta_j(\overline{\mathbf{f}})\right)$ satisfy the strict ordering
\[
1 = \hat{Q}_1 > \hat{Q}_2 > \cdots > \hat{Q}_{K} > 0,
\]
with the boundary condition $\hat{Q}_{K+1} \triangleq 0$.
\end{theorem}

\textbf{Proof.}
Since $\sigma(\cdot)$ maps every finite logit to the open interval $(0,1)$, each conditional factor satisfies $\hat{p}_k = \sigma(\theta_k(\overline{\mathbf{f}})) \in (0,1)$ strictly. By the recurrence $\hat{Q}_{k+1} = \hat{Q}_k \cdot \hat{p}_k$, it follows that $\hat{Q}_{k+1} < \hat{Q}_k$ for all $k$, establishing strict monotone decrease from $\hat{Q}_1 = 1$ downward. Since every factor is strictly positive, the cumulative product remains positive, ensuring $\hat{Q}_K > 0$. The terminal condition $\hat{Q}_{K+1} = 0$ is an architectural constraint arising from the absence of a $K$-th threshold logit, not a learned parameter. Consequently, this factorization structurally precludes the non-monotone constraint violations endemic to standard ordinal regression models, ensuring robust behavior even across the severely imbalanced grades of the LSMA-PQR dataset.

The continuous disc height estimation is optimized via a Smooth L1 regression loss, defined as $\mathcal{L}_\text{ht} = \text{SmoothL1}(\hat{h}, h)$. Concurrently, the multi-label disorder classification is supervised using a standard BCE loss, $\mathcal{L}_\text{dis} = -\frac{1}{9}\sum_{j=1}^{9}[d_j \log \hat{d}_j + (1{-}d_j)\log(1{-}\hat{d}_j)]$, across the nine target pathologies.


By aggregating these task-specific losses, the shared visual manifold is subjected to competing gradient updates; rather than degrading performance, this composite optimization \eqref{eq:loss} acts as a crucial implicit regularizer that enriches the shared feature representation beyond any single-task baseline, thereby providing theoretical grounding for the multi-task gains observed empirically (\Cref{sec:ablation_results}).

\begin{theorem}
\label{thm:multitask}
Let $\mathcal{L}_{\mathrm{MT}}(\mathbf{f}) = \mathcal{L}_{\mathrm{seg}}(\mathbf{f}) + \mathcal{L}_{\mathrm{grade}}(\mathbf{f})$ denote the joint objective over shared features $\mathbf{f}$. Any feature direction $\mathbf{v}$ satisfying $\nabla_{\mathbf{v}}\mathcal{L}_{\mathrm{seg}} = 0$ but $\nabla_{\mathbf{v}}\mathcal{L}_{\mathrm{grade}} \neq 0$ is preserved under $\mathcal{L}_{\mathrm{MT}}$ but discarded under single-task segmentation optimization.
\end{theorem}

\textbf{Proof.}
Under $\mathcal{L}_{\mathrm{seg}}$ alone, the gradient update along $\mathbf{v}$ is $\Delta\mathbf{f}_{\mathbf{v}} \propto -\nabla_{\mathbf{v}}\mathcal{L}_{\mathrm{seg}} = 0$. Since the update vanishes, $\mathbf{v}$ lies in the null space of the segmentation objective; it is not reinforced and may be suppressed by weight decay. Under $\mathcal{L}_{\mathrm{MT}}$, gradient superposition yields
\[
\Delta\mathbf{f}_{\mathbf{v}} \propto -\bigl(\nabla_{\mathbf{v}}\mathcal{L}_{\mathrm{seg}} + \nabla_{\mathbf{v}}\mathcal{L}_{\mathrm{grade}}\bigr) = -\nabla_{\mathbf{v}}\mathcal{L}_{\mathrm{grade}} \neq 0,
\]
actively updating and preserving $\mathbf{v}$. Directions such as $\mathbf{v}$ encode non-spatial clinical biomarkers (\eg, T2-weighted signal intensity diagnostic of disc degeneration) that carry grading information but do not define spatial boundaries. Joint training therefore acts as a representation regularizer, retaining clinically informative feature dimensions beyond what segmentation supervision alone would preserve \cite{caruana1997multitask}.

Together, this unified composite objective ensures that the shared latent representations encode both fine-grained spatial geometries and high-level clinical semantics.

\subsection{Report Generator}
\label{sec:report_gen}
Structured radiological reports are generated from model predictions using a two-stage retrieval-augmented template filling (RATF) pipeline, as depicted in \Cref{fig:ratf}. In the first stage, template filling, per-level predictions including the Pfirrmann grade, disorder probabilities exceeding $0.5$, and disc height estimates are deterministically mapped to predefined sentence templates aligned with the American College of Radiology (ACR) Lumbar Spine Reporting Guidelines. Each analyzed disc level, from L3-L4 through L5-S1, generates an independent finding sentence. The second stage, retrieval-augmented template filling (RATF), employs a FAISS index \cite{johnson2019billion} containing reference report sentences curated from radiologist-authored reports within the training cohort. This index is queried using the all-MiniLM-L6-v2 \cite{wang2020minilm} embedding of the initially generated template sentence. 

\begin{figure}[!ht]
	\centering
	\includegraphics[width=\columnwidth]{figure3_ratf.pdf}
	\caption{Architecture of the retrieval-augmented template filling (RATF) report generator.}
	\label{fig:ratf}
\end{figure}

To enrich the phrasing and embed contextual clinical detail, the top retrieved sentences, ranked by cosine similarity, are evaluated. If the highest similarity score exceeds a predefined threshold ($\tau=0.75$), the retrieved sentence replaces the template sentence; otherwise, the original template is retained verbatim. This mechanism ensures that retrieval augmentation enhances linguistic diversity without degrading the clinical accuracy of the deterministic baseline. The resulting enriched, per-level findings are subsequently rendered as concatenated natural-language sentences with explicit disc-level context (\eg, ``At L4--L5, a disc bulge is noted with mild thecal sac compression''). This approach produces flowing radiological prose that improves lexical alignment with radiologist-authored free-text reports while preserving complete clinical content fidelity.

% \emph{Note:} RATF is a template-based system with retrieval-augmented phrasing enrichment; it does not constitute an end-to-end generative language model and is subject to the limitations discussed in \Cref{sec:limitations}.



% ══════════════════════════════════════════════════════════════════
\section{Experimental Setup}
\label{sec:experiments}

% \subsection{Dataset}
% \label{sec:dataset}



We use the LSMA-PQR dataset \cite{masood_lsma_pqr_2025}, comprising of $N=515$ consecutive symptomatic patients, each contributing eight multi-planar acquisitions including three axial disc-level slices (L3-L4, L4-L5, L5-S1) in both T1- and T2-weighted sequences and two mid-sagittal slices (T1w, T2w), yielding 4,120 images in total (3,090 axial, 1,030 sagittal) at a standardized $384{\times}384$ pixel resolution. The annotation framework encompasses four complementary modalities:
\begin{enumerate}[label=(\roman*)]
    \item Pixel-level segmentation masks for seven anatomical structures, including three axial classes (intervertebral disc, posterior elements and spinal canal) and four sagittal classes (vertebral body, IVD, sacrum and spinal canal), provided in YOLO-polygon and indexed-mask formats with verified cross-format spatial consistency (mean IoU $= 0.994 \pm 0.008$ across all 4,120 images), totalling 25,143 annotation instances.
    \item Pfirrmann degeneration grades (I--V) for 1,545 disc-level assessments ($515 \times 3$ levels), on mid-sagittal T2-weighted images.
    \item Intervertebral disc heights measured from mid-sagittal T2-weighted images.
    \item 515 structured radiological reports derived from free-text clinical documents and binary disorder labels for 9 pathological findings.
\end{enumerate}

Dataset characteristics are summarized in \Cref{tab:dataset}. The dataset is partitioned into patient-level stratified splits to prevent data leakage: 360 patients (2,880 images) for training, 77 patients (616 images) for validation and 78 patients (624 images) for testing (70/15/15\%). Distribution is based on the maximum Pfirrmann grade per patient to maintain proportional grade representation across all splits.

\begin{table}[htbp]
\centering
\caption{LSMA-PQR dataset characteristics by split. Grading labels are assigned per disc level on mid-sagittal T2-weighted images and propagated to all imaging planes as multi-view training supervision; the ``Graded discs'' row counts axial images only, each corresponding to one disc level. Grade distribution is shown for train/val/test splits to verify stratification.}
\label{tab:dataset}
\begin{tabular}{lccc}
\toprule
& Train & Val & Test \\
\midrule
Patients & 360 & 77 & 78 \\
Images & 2,880 & 616 & 624 \\
Axial / Sagittal & 2,160 / 720 & 462 / 154 & 468 / 156 \\
Graded discs & 2,160 & 462 & 468 \\
\midrule
Grade I & 36 & 8 & 2 \\
Grade II & 1,164 & 264 & 238 \\
Grade III & 514 & 108 & 130 \\
Grade IV & 386 & 72 & 90 \\
Grade V & 60 & 10 & 8 \\
\bottomrule
\end{tabular}
\end{table}

\subsection{Implementation Details}
\label{sec:implementation}

All experiments use PyTorch 2.5.1 with mixed-precision (BFloat16) training on a single NVIDIA RTX 4070 GPU (8 GB VRAM). Images are resized to $384 \times 384$ pixels with standard ImageNet normalization. Data augmentation via Albumentations \cite{buslaev2020albumentations} includes horizontal flips ($p=0.5$), shift-scale-rotate (shift $\pm 0.05$, scale $\pm 0.1$, rotate $\pm 15^\circ$), Gaussian blur ($\sigma \in [0.1, 2.0]$) and random brightness--contrast adjustments. Elastic deformation, a standard augmentation for deformable anatomy, was omitted after preliminary experiments showed negligible benefit, likely due to the relatively rigid geometry of lumbar vertebral structures compared to soft-tissue organs.
The training protocol follows a two-phase schedule: Phase 1 (5 epochs) uses a learning rate of $10^{-4}$ with AdamW ($\beta_1=0.9$, $\beta_2=0.999$, weight decay $10^{-4}$) and batch size 2 with gradient accumulation over 8 steps (effective batch size 16); Phase 2 (up to 50 epochs) reduces the learning rate to $5{\times}10^{-5}$ and applies cosine annealing with warm restarts ($T_0=10$, $T_\text{mult}=2$).
Gradient clipping at norm 1.0 and automatic mixed precision with dynamic loss scaling are applied throughout both phases. Model selection uses the best mean Dice on the validation set, with early stopping (patience 10 epochs) in Phase 2 only. Class-frequency inverse weights are precomputed from training-set pixel counts and applied as \texttt{pos\_weight} in the per-class segmentation BCE loss, following $w_c = N_\text{total} / (C \cdot N_c)$, where $N_c$ is the class pixel count. This rebalancing is critical for minority classes such as the axial spinal canal ($w_\text{SC}=90.0$; $<$0.25\% of pixels).

\subsection{Ablation Design}
\label{sec:ablation}
To isolate the contribution of each architectural component, we conduct a systematic ablation study comprising four variants of increasing complexity:
\begin{itemize}
    \item \textbf{A0}: No LoRA, no Prompt Calibrator (1.00 M trainable params)
    \item \textbf{A1}: LoRA only, no Prompt Calibrator (2.19 M params)
    \item \textbf{A2}: LoRA + Prompt Calibrator (2.73 M params), the \emph{full model}
    \item \textbf{A3}: Full model + increased boundary loss ($\lambda_\text{bound}=1.0$; 2.73 M params)
\end{itemize}

All four variants undergo the identical two-phase training protocol described above, ensuring that observed performance differences reflect purely architectural contributions without confounding by training duration or hyperparameter variation. 

\subsection{Baseline Comparison}
We evaluate \modelname{} against two categories of baselines.
\emph{Direct baselines} are evaluated on the LSMA-PQR test set under identical conditions: SAM ViT-B with automatic mask generation \cite{kirillov2023segment} (no prompts; generated masks matched to ground truth via IoU-based assignment), SAM ViT-B with bounding boxes derived from mask annotations (representing an upper bound for prompt-based approaches), and MedSAM2 \cite{medsam2_2024} with Hiera-Tiny backbone and bounding-box prompts.
Additionally, nine foundation-LLM report baselines including MedGemma-4B-IT \cite{medgemma2025}, BioMistral-7B \cite{biomistral2024}, Meditron3-8B \cite{sallinen2025llama3meditron}, BioGPT-Large \cite{luo2022biogpt}, LLaVA-Rad-7B \cite{zambranochaves2025llavarad}, Phi-3.5-mini \cite{phi35mini2024}, Llama-3.2-3B-Instruct \cite{llama32_2024}, SmolLM2-1.7B-Instruct \cite{smollm2_2024} and Qwen2.5-1.5B-Instruct \cite{qwen25_2024} are evaluated for report-generation comparison against RATF under a matched-patient protocol (\Cref{sec:report_results}).

To evaluate segmentation performance within the broader medical imaging landscape, we further compare against \emph{literature baselines} comprising of published results from established methods on related spine and multi-organ benchmarks. Among task-specific architectures, nnU-Net \cite{isensee2021nnu}, a self-configuring segmentation framework, achieves $\bar{D}=0.90$--$0.93$ on the SPIDER benchmark \cite{van2024spider} for vertebral body and disc segmentation, while SwinUNETR \cite{hatamizadeh2022swinunetr}, a Swin Transformer-based encoder-decoder, reports $\bar{D}=0.87$--$0.91$ on 3D volumetric spine tasks. Among foundation models adapted to medical imaging, SAM-Med2D \cite{cheng2023sammed2d} (SAM fine-tuned on 4.6M medical images with point and box prompts) attains $\bar{D}=0.82$--$0.86$ on musculoskeletal structures, while MedSAM \cite{ma2024medsam} reaches $\bar{D}=0.85$--$0.89$ on abdominal and cardiac benchmarks with limited spine-specific evaluation. Prompt-driven models are represented by SAT-Pro \cite{satpro2025} (497-category text-prompted 3D segmentation; $\bar{D}=0.85$--$0.89$) and MedUniSeg \cite{meduniseg_2024} (prompt-driven universal segmentation; $\bar{D}=0.80$--$0.85$ on multi-organ benchmarks), while DAPSAM \cite{dapsam_2025} introduces domain-adaptive prototypical prompts ($\bar{D}=0.83$--$0.87$ on cross-domain medical segmentation).
These literature values are not directly comparable to our results due to differences in datasets, annotation granularity and evaluation protocols; nevertheless, they provide a valuable reference frame for interpreting \modelname{}'s performance. \Cref{fig:forest_baselines} presents a forest-plot summary of these literature benchmarks alongside \modelname{}.
With the experimental framework established, the following section presents the quantitative and qualitative outcomes across all evaluation axes.

% ======================================================================
\section{Results}
\label{sec:results}
In this section, we present a comprehensive evaluation of the \modelname{} framework. We begin by isolating the contributions of key architectural components through an ablation study and comparing our approach against established baselines. We then detail the model's performance across its multi-task clinical objectives: semantic segmentation, Pfirrmann grading and disorder detection. Next, we evaluate the quality of the structured radiology reports synthesized by our Retrieval-Augmented Template Filling (RATF) pipeline against foundation large language models. Finally, we assess the model's generalizability through zero-shot cross-dataset validation.

\subsection{Ablation Study}
\label{sec:ablation_results}

Visualization of DSC and validation IoU training curves across all four ablation configurations is shown in \Cref{fig:training_curves}. \emph{A0} (gray) serves as the frozen-backbone baseline with no domain adaptation or prompt conditioning. It converges rapidly within the warmup phase but plateaus at a low DSC ${\approx} 0.40$ and IoU ${\approx} 0.26$, remaining flat throughout all 50 epochs and indicating that the frozen DINOv2 features alone are insufficient for lumbar spine segmentation. Training was interrupted before convergence (diamond marker), confirming no further gains were expected. \emph{A1} (blue) adds LoRA fine-tuning to A0. Performance remains similarly stagnant at DSC ${\approx} 0.40$ and IoU ${\approx} 0.27$, suggesting that low-rank weight updates alone, without domain-specific adaptation or prompt conditioning, yield negligible improvement over the frozen baseline. A1 was also interrupted early. \emph{A2} (green) introduces both LoRA fine-tuning and the Prompt Calibrator, enabling vocabulary-conditioned feature gating of the DINOv2 backbone. This configuration shows the strongest learning dynamics, with DSC rising steeply from epoch 5 onward and reaching ${\approx} 0.81$ at best validation (star marker, epoch ${\approx} 50$), and IoU reaching ${\approx} 0.70$. The train--val DSC gap remains modest, indicating stable generalisation without overfitting. \emph{A3} (red) extends A2 by increasing the boundary loss weighting to $\lambda_\text{bound}=1.0$, placing greater emphasis on anatomical edge delineation. It achieves the highest DSC of ${\approx} 0.87$ and IoU of ${\approx} 0.78$ at epoch 50 (star markers), consistently outperforming all ablated variants from epoch 5 onward. The training DSC (dashed red) closely tracks the validation curve, suggesting stable optimisation without overfitting. GPU memory consumption (right $y$-axis, grey scatter) remains well below the 8 GB budget across all configurations, confirming the memory efficiency of the proposed design. The warmup boundary (vertical dotted line at epoch 5) is visible throughout.

\begin{figure}[!htbp]
    \centering
    \includegraphics[width=\columnwidth]{figure4.png}
    \caption{Training dynamics across ablation variants (test-set Dice and validation IoU over 55 epochs). All variants follow the identical two-phase protocol (5 + 50 epochs). Stars mark the best-validation-Dice epoch per variant; diamonds indicate early-stopped runs. A2 (full model) rises steeply from epoch 5 and continues improving throughout Phase 2; A0 and A1 converge early due to limited model capacity. Dotted vertical line: warmup boundary at epoch 5. GPU peak VRAM (grey scatter, right axis) remains well below the 8 GB budget (dashed red)}.
    \label{fig:training_curves}
\end{figure}

The same ablation scheme was evaluated on the test set ($n=624$ images from 78 patients) as shown in \Cref{tab:ablation} and \Cref{fig:ablation}. The full model (A2) achieves the best multi-task balance with mean Dice $\bar{D}=0.861$ (95\% CI $[0.852, 0.869]$) and mean IoU $\bar{J}=0.770$ (95\% CI $[0.758, 0.781]$), computed via 10,000-iteration bootstrap resampling of per-image scores. A3 ($\lambda_{\text{bound}}=1.0$) achieves marginally higher segmentation ($\bar{D}=0.871$, $p=0.003$) but lower grade accuracy (0.583 as compared to 0.618). All variants underwent the identical two-phase training protocol (5 + up to 50 epochs with early stopping), ensuring that performance differences reflect architectural contributions.

\begin{table}[!htbp]
\centering
\caption{Ablation study results on the test set ($n=624$). $\bar{D}$: mean prompted Dice; $\bar{J}$: mean prompted IoU; 95\% bootstrap CIs in brackets. All variants trained with the identical two-phase protocol (Phase 1: 5 epochs; Phase 2: up to 50 epochs with early stopping patience 10). Best values in \textbf{bold}.}
\label{tab:ablation}
\resizebox{\columnwidth}{!}{
\begin{tabular}{lcccccc}
\toprule
Variant & Params & $\bar{D}\uparrow$ & $\bar{J}\uparrow$ & Gr. Acc.$\uparrow$ & Ht. MAE$\downarrow$ & VRAM \\
\midrule
A0 & 1.00M & 0.378 [.36,.40] & 0.261 [.25,.28] & 0.528 & 1.71 & 0.54 GB \\
A1 & 2.19M & 0.375 [.36,.39] & 0.260 [.24,.28] & \textbf{0.652} & 1.61 & 0.62 GB \\
A2 & 2.73M & 0.861 [.85,.87] & 0.770 [.76,.78] & 0.618 & \textbf{1.56} & 0.63 GB \\
A3 & 2.73M & \textbf{0.871} [.86,.88] & \textbf{0.785} [.77,.80] & 0.583 & 1.61 & 0.63 GB \\
\bottomrule
\end{tabular}
}
\end{table}

\begin{figure}[!htbp]
    \centering
    \includegraphics[width=\columnwidth]{figure5.pdf}
    \caption{Ablation study: mean Dice ($\bar{D}$), mean IoU ($\bar{J}$) and grade accuracy across ablation variants on the test set.}
    \label{fig:ablation}
\end{figure}

% All pairwise comparisons (Wilcoxon signed-rank test\footnote{Wilcoxon signed-rank tests are applied to paired per-image Dice scores across ablation variants, with Bonferroni correction for multiple comparisons ($\alpha = 0.05/6$).}, Bonferroni-corrected) confirm statistically significant differences: A3 vs. A2: $W=84,074$, $p=2.9 \times 10^{-3}$; A3 vs. A0: $W=689$, $p=2.0 \times 10^{-102}$; A3 vs. A1: $W=634$, $p=1.5 \times 10^{-102}$.

\subsection{Per-Class Segmentation}
Having established the overall superiority of the full model configuration, we next examine how performance is distributed across individual anatomical structures. \Cref{tab:perclass} reports per-class Dice scores for all ablation variants. A2 achieves strong performance on IVD (0.948), VB (0.899), Sacrum (0.892) and SC\textsubscript{sag} (0.886). The axial spinal canal (SC) is the most challenging structure due to its extremely small cross-sectional area ($\sim$365 pixels per $384 \times 384$ image, $<$0.25\% of total pixels).
Application of class-balanced BCE with inverse-frequency pos\_weight ($w_\text{SC}=90.0$) is essential to prevent the dedicated segmentation channel from collapsing under severe class imbalance; without this rebalancing, the axial SC Dice degrades to near-zero. The prompted segmentation head achieves $D=0.849$ when conditioned on ``spinal canal'' text prompts (149 test images), confirming that the DINOv2 features contain sufficient SC discriminability; the bottleneck was in the unweighted per-channel BCE loss, not the feature representation. The mean foreground Dice across all 7 classes is $\bar{D}_{\text{fg}} = 0.788$. At lumbar levels (L3--S1), the ``SC'' label corresponds to the thecal sac containing the cauda equina rather than the spinal cord proper, which terminates at the conus medullaris near L1-L2.

\begin{table}[!htbp]
\centering
\caption{Per-class Dice scores ($D$) on the test set. IVD: intervertebral disc, PE: posterior elements (laminae, pedicles, spinous processes, facet joints), SC: spinal canal, VB: vertebral body. Subscript ``sag'' denotes sagittal variants. $\bar{D}_{\text{fg}}$: mean per-image foreground Dice (each image contributes its mean Dice across applicable foreground classes, then averaged over images); $\bar{D}_{\text{fg}}^{*}$: same computation excluding axial SC.}
\label{tab:perclass}
\resizebox{\columnwidth}{!}{
\begin{tabular}{lccccccccc}
\toprule
 & IVD & PE & SC & VB & IVD\textsubscript{sag} & Sacrum & SC\textsubscript{sag} & $\bar{D}_{\text{fg}}$ & $\bar{D}_{\text{fg}}^{*}$ \\
\midrule
A0 & 0.910 & 0.634 & 0.217 & 0.670 & 0.475 & 0.773 & 0.763 & 0.608 & 0.704 \\
A1 & 0.881 & 0.715 & 0.382 & 0.742 & 0.554 & 0.729 & 0.825 & 0.673 & 0.721 \\
A2 & 0.948 & 0.839 & 0.491 & \textbf{0.899} & 0.824 & 0.892 & \textbf{0.886} & 0.788 & \textbf{0.881} \\
A3 & \textbf{0.948} & \textbf{0.853} & \textbf{0.527} & 0.848 & \textbf{0.850} & \textbf{0.901} & 0.884 & \textbf{0.800} & 0.881 \\
\bottomrule
\end{tabular}
}
\end{table}

Consistent performance across three lumbar disc is observed, shown in \Cref{fig:per_level}, with no systematic degradation at specific anatomical levels. The marginally lower mean and wider lower tail at L5-S1 are consistent with increased anatomical variability at the lumbosacral junction.

\begin{figure}[!htbp]
    \centering
    \includegraphics[width=0.9\columnwidth]{figure6.pdf}
    \caption{Per-disc-level prompted Dice distributions (A2, axial test set, $n=468$). Mean Dice is consistently high at L3-L4 ($0.866$), L4-L5 ($0.859$) and L5-S1 ($0.830$).}
    \label{fig:per_level}
\end{figure}

% (SC Dice$=$0.491 for A2, \Cref{tab:perclass})
% (\Cref{sec:limitations}, item 2)



% \subsubsection{Axial Spinal Canal Enhancement}
% \label{sec:sc_enhancement}
To address the axial spinal canal segmentation challenge, we developed a dedicated SC-enhanced model using coarse-to-fine segmentation and targeted small-structure augmentation. As observed in the baseline evaluations (\Cref{tab:perclass}), the extreme class imbalance severely suppressed SC segmentation performance, yielding baseline Dice scores of $0.491$ and $0.527$ for the A2 and A3 variants, respectively.

The enhancement protocol employs three mechanisms. First, SC-focused zoom augmentation uses a fifty percent probability of cropping a region between $128$--$256$ pixels centered on the SC centroid and resizing it to the native $384{\times}384$ resolution. Second, elastic deformation and grid distortion augmentations are applied to improve robustness to morphological variations. Third, selective SC-weighted loss multipliers are utilized. For the A2 variant, an aggressive weighting scheme ($w_\text{SC}=10{\times}$) was applied. In contrast, the A3 variant, which natively benefits from boundary-aware regularization, utilized a more conservative weighting ($w_\text{SC}=2.5{\times}$) to preserve the structural integrity of adjacent anatomies.

\begin{table}[htbp]
\centering
\caption{Impact of the targeted axial spinal canal (SC) enhancement protocol evaluated on the axial test set ($n=468$). $\Delta$ denotes the absolute Dice improvement, demonstrating the improvement of minority-class performance ($\text{SC}$) without degrading the spatial performance of adjacent structures ($\text{IVD}$, $\text{PE}$). Best post-enhancement values per variant are highlighted in \textbf{bold}.}
\label{tab:sc_enhanced_combined}
\resizebox{\columnwidth}{!}{%
\begin{tabular}{@{}l ccc c ccc@{}}
\toprule
& \multicolumn{3}{c}{\textbf{A2} ($w_\text{SC}=10{\times}$)} & \phantom{a} & \multicolumn{3}{c}{\textbf{A3} ($w_\text{SC}=2.5{\times}$)} \\
\cmidrule{2-4} \cmidrule{6-8}
Anatomy & Baseline & Enhanced & $\Delta$ && Baseline & Enhanced & $\Delta$ \\
\midrule
IVD & 0.948 & \textbf{0.970} & +0.023 && 0.948 & 0.955 & +0.007 \\
PE  & 0.839 & \textbf{0.915} & +0.076 && 0.853 & 0.871 & +0.018 \\
SC  & 0.491 & \textbf{0.910} & \textbf{+0.418} && 0.527 & 0.787 & +0.260 \\
\midrule
$\bar{D}_\text{ax}$ & 0.759 & \textbf{0.931} & +0.173 && 0.776 & 0.871 & +0.095 \\
\bottomrule
\end{tabular}
}
\end{table}

The quantitative impact of this strategy is given in \Cref{tab:sc_enhanced_combined}. For the A2 variant, the enhancement strategy increases the SC Dice from $0.491$ to $0.910$ ($+85\%$ improvement). For the A3 variant, the targeted boost increases the SC Dice from $0.527$ to $0.787$ ($+49\%$ relative improvement). These gains confirm that the shared DINOv2 feature representations contain sufficient discriminative capacity to isolate small structures and that the initial low performance was strictly an artifact of extreme pixel minority, not an encoder bottleneck. 


The per-level performance improvement from A3 baseline to SC-Enhanced model is plotted in \Cref{fig:sc_enhancement_dumbbell}. This visualization maps the absolute performance delta ($\Delta$), confirming that targeted optimization for the minority SC class yields a maximum gain of $+0.301$ at L5-S1 without causing degradation in adjacent structures. Specifically, Posterior Elements (PE) experience a marginal regularization benefit ($\Delta \approx +0.01$ to $+0.02$), while Intervertebral Disc (IVD) spatial score remains highly stable near $0.95$. These results visually reinforce that the coarse-to-fine anatomical zoom resolves minority-class collapse while preserving the integrity of the shared multi-task latent space.

\begin{figure}[htbp]
    \centering
    \includegraphics[width=0.9\columnwidth]{figure7.pdf}
    \caption{%
        Per-level prompted Dice score improvements under the SC-Enhanced training protocol. Circles and diamonds denote the A3 baseline and the optimized SC-Enhanced model, respectively.}
    \label{fig:sc_enhancement_dumbbell}
\end{figure}



\subsection{Direct and Literature Baselines Comparison}
\label{sec:baseline_comparison}

With the architectural contributions of each component established, we now situate \modelname{} against existing segmentation methods. \Cref{tab:baselines} compares \modelname{} against direct baselines evaluated on LSMA-PQR and literature baselines from related spine benchmarks.

\begin{table*}[t]
\centering
\caption{Comparison of \modelname{} against evaluated methods. Direct baselines are evaluated on the LSMA-PQR test set ($n=624$); literature values ($^*$) are from published results on related benchmarks and are included for contextual reference only. ``GT BBox'' indicates the requirement of ground-truth bounding-box annotations at inference; ``Text Prompt'' indicates zero-shot text conditioning. $t_{\text{inf}}$: per-image inference time. \modelname{} is the only framework unifying text-guided prompting, multi-task clinical outputs and extreme parameter efficiency.}
\label{tab:baselines}
% \tabular* with \extracolsep{\fill} forces the table to perfectly span the \textwidth
\begin{tabular*}{\textwidth}{@{\extracolsep{\fill}} l c c c c c c c c @{}}
\toprule
Method & $\bar{D}\uparrow$ & $\bar{J}\uparrow$ & Params & GT BBox & Text Prompt & Multi-task & $t_{\text{inf}}$ & Dataset \\
\midrule
\multicolumn{9}{@{}l}{\textit{Direct baselines (evaluated on LSMA-PQR test set)}} \\
\addlinespace
SAM ViT-B Auto \cite{kirillov2023segment} & 0.433 & 0.366 & 93.7M & No & No & No & 3.49 s & Ours \\
MedSAM2 \cite{medsam2_2024} & 0.755 & 0.654 & 38.9M & Yes & No & No & 0.011 s & Ours \\
SAM ViT-B \cite{kirillov2023segment} & 0.810 & 0.711 & 93.7M & Yes & No & No & 0.15 s & Ours \\
\textbf{\modelname{}} & \textbf{0.861} & \textbf{0.770} & \textbf{2.73M} & \textbf{No} & \textbf{Yes} & \textbf{Yes} & \textbf{0.039 s} & Ours \\
\midrule
\multicolumn{9}{@{}l}{\textit{Literature baselines (on different datasets)}} \\
\addlinespace
nnU-Net \cite{isensee2021nnu} & 0.90--0.93 & -- & $\sim$31M & No & No & No & -- & SPIDER \\
SwinUNETR \cite{hatamizadeh2022swinunetr} & 0.87--0.91 & -- & 62M & No & No & No & -- & BraTS \\
SAM-Med2D \cite{cheng2023sammed2d} & 0.82--0.86 & -- & 93.7M & Yes & No & No & -- & SA-Med2D \\
MedSAM \cite{ma2024medsam} & 0.85--0.89 & -- & 93.7M & Yes & No & No & -- & Multi \\
SAT-Pro \cite{satpro2025} & 0.85--0.89 & -- & 104M & No & Yes & No & -- & Multi \\
MedUniSeg \cite{meduniseg_2024} & 0.80--0.85 & -- & 93.7M & Yes & Yes & No & -- & Multi \\
DAPSAM \cite{dapsam_2025} & 0.83--0.87 & -- & 93.7M & Yes & No & No & -- & Multi \\
\bottomrule
\end{tabular*}
\end{table*}


As exhibited in the \Cref{tab:baselines}, \modelname{} achieves a higher mean Dice similarity coefficient than all direct baselines evaluated on the LSMA-PQR dataset. Specifically, our framework performed 6.3\% better than SAM ViT-B utilizing ground-truth bounding boxes. Furthermore, MedSAM2 exhibited a 12.3\% degraded performance compared to our model, despite both baselines relying on exact spatial annotations, while our framework operated zero-shot via text prompts using only 2.9\% of the trainable parameters. This performance gap demonstrates that the learned vocabulary conditioning of the Prompt Calibrator is fundamentally more effective than spatial bounding-box prompting for dense anatomical segmentation in this domain. Additionally, SAM ViT-B relying solely on automatic mask generation showed a 49.7\% degraded performance compared to our model, confirming that unprompted, general-purpose segmentation frameworks are insufficient for clinical anatomy. When contextualized against literature baselines, the text-prompted capability of \modelname{} is highly competitive with SAM-Med2D and SAT-Pro and closely approaches MedSAM, despite utilizing substantially fewer parameters and lacking spatial prompting. Finally, while dedicated single-task architectures like nnU-Net establish the highest reported performance on related benchmarks—performing roughly 6.8\% better than our text-guided framework—they strictly require task-specific architecture search and full-parameter model training for every distinct clinical domain, whereas \modelname{} unifies these objectives in a single lightweight forward pass. \Cref{fig:forest_baselines} contextualises these literature benchmarks against \modelname{}'s bootstrap-estimated performance range.

\begin{figure}[!htbp]
    \centering
    \includegraphics[width=\columnwidth]{figure8.pdf}
    \caption{Forest plot comparing \modelname{} against literature baselines on Mean Dice Similarity Coefficient ($\bar{D}$) for reference purposes. Blue bars denote published DSC ranges on respective benchmarks (SPIDER, BraTS, SA-Med2D, multi-organ); the red bar indicates \modelname{}'s 95\% bootstrap CI on LSMA-PQR ($n=624$).}
    \label{fig:forest_baselines}
\end{figure}


%%%%%%%%%% COMPARISON WITH DIRECTLY TRAINED BASELINES, SEGMENTATION ONLY TASK %%%%%%%%%%
% {\color{red}
% \subsubsection{Direct Baseline Comparison on LSMA-PQR}
% \label{sec:direct_baselines}
% To address the limitation of relying solely on literature-reported baselines (\Cref{sec:limitations}), we train three representative segmentation architectures directly on LSMA-PQR using the identical train/val/test splits, $384{\times}384$ input resolution and per-class Dice/IoU evaluation protocol.
% All baselines are segmentation-only models (no grading, height, disorder, or report generation).

% \begin{table}[htbp]
% \color{red}
% \centering
% \caption{Direct baseline comparison on LSMA-PQR test set ($n=624$). All models trained from scratch on the same splits with mixed-precision (bf16) training. \modelname{} is the only multi-task model; all others are segmentation-only.}
% \label{tab:direct_baselines}
% \resizebox{\columnwidth}{!}{
% \begin{tabular}{lccccccc}
% \toprule
% Method & Params & $\bar{D}\!\uparrow$ & IVD & PE & SC & VB & SC\textsubscript{sag} \\
% \midrule
% nnU-Net (PlainConvUNet) & 31.2M & \textbf{0.944} & \textbf{0.977} & \textbf{0.932} & \textbf{0.919} & \textbf{0.962} & \textbf{0.933} \\
% SwinUNETR (2D) & 6.3M & 0.937 & 0.975 & 0.925 & 0.903 & 0.958 & 0.925 \\
% SAM ViT-B Fine-Tuned & 90.3M\textsuperscript{\dag} & 0.914 & 0.966 & 0.897 & 0.869 & 0.947 & 0.912 \\
% \textbf{\modelname{} (A2)} & \textbf{2.7M} & 0.788 & 0.948 & 0.839 & 0.491 & 0.899 & 0.886 \\
% \bottomrule
% \multicolumn{8}{l}{\footnotesize \textsuperscript{\dag}90.3M total; only 0.66M decoder parameters are trainable (encoder frozen).}
% \end{tabular}
% }
% \end{table}

% As expected, dedicated segmentation architectures outperform the multi-task \modelname{} on raw foreground Dice: nnU-Net achieves $\bar{D}=0.944$, SwinUNETR $0.937$ and SAM ViT-B (fine-tuned decoder) $0.914$, versus \modelname{}'s $0.788$.
% The performance gap is most pronounced for the axial spinal canal (SC Dice: $0.919$ vs. $0.491$), where dedicated models benefit from full-resolution feature maps.
% However, this comparison highlights a fundamental architectural trade-off: \modelname{} simultaneously produces five clinical outputs (segmentation, grading, height, disorders, reports) using $11.5{\times}$ fewer trainable parameters than nnU-Net, which requires a separate model for each clinical task.
% The single-task baselines require $N{\times}$ the total parameter budget for $N$ clinical tasks, whereas \modelname{} achieves all five in a single 39 ms forward pass.
% } % end \color{red}

To evaluate computational efficiency, all models were benchmarked on a single RTX 4070 Laptop GPU (8 GB VRAM) using uniform configuration. \modelname{} achieves the highest segmentation performance ($\bar{D}=0.861$) while optimizing only 2.73M trainable parameters (3.08\% of its 88.8M total), requiring 0.63 GB of peak VRAM and executing a single forward pass in 39 ms. In contrast, SAM ViT-B ($\bar{D}=0.810$) and SAM Auto ($\bar{D}=0.433$) train the full 93.7M parameter space, consume over 4 GB of VRAM and exhibit inference times of 150 ms and 3.49 s, respectively. While MedSAM2 provides the fastest inference (11 ms) and requires 1.8 GB VRAM, it yields lower segmentation performance ($\bar{D}=0.755$) despite utilizing 38.9M trainable parameters. Finally, the MedGemma foundation model (4.97B parameters), evaluated exclusively for report generation (\Cref{sec:report_results}), exceeds the 8 GB VRAM budget, necessitating CPU offloading and resulting in inference times approaching 200 s.


\subsection{Pfirrmann Grading}
\label{sec:grading_results}

Beyond segmentation, \modelname{}'s multi-task head yields ordinal Pfirrmann grading without additional annotations. On the test set ($n=468$ graded disc levels), A2 achieves classification accuracy of 61.8\% with 92.3\% of predictions within $\pm 1$ grade and a grade MAE of 0.46. To validate these results within the clinical domain, Pfirrmann grades were assessed by an expert radiologist on mid-sagittal T2-weighted images following the original protocol \cite{pfirrmann2001magnetic}; the resulting per-disc-level grade labels are propagated to all imaging planes (axial T1/T2, sagittal T1/T2) of the same disc level as a multi-view training signal. Due to the natural prevalence of mild degeneration in the clinical screening, Grade II constitutes the majority class, comprising $50.9\%$ ($238/468$) of the test samples. Consequently, a baseline naive classifier defaulting to Grade II would inherently achieve an exact accuracy of $50.9\%$ and a within-$\pm 1$ accuracy of $79.1\%$, encompassing Grades I, II and III. \modelname{} substantially exceeds these baselines by $+10.9$ and $+13.2$ percentage points, respectively, confirming that the architecture learns genuine discriminative morphological features rather than merely exploiting class prevalence.Inter-observer agreement among radiologists for Pfirrmann grading is reported at $\kappa = 0.56$--$0.76$ \cite{wilke2006discdegeneration}, corresponding to approximately 70--80\% exact agreement, placing our model's 61.8\% accuracy within the range of human variability, though below the upper bound.

Per-grade analysis (\Cref{fig:confusion}) reveals strong Grade II accuracy (87.0\%, F1$=$0.758) but reduced performance on rare extreme grades: Grade I ($n=2$) is not predicted correctly as it has only 2 samples which are insufficient for statistically reliable per-class metric estimation; Grade V ($n=8$) achieves only 25\% recall (\Cref{fig:grade_f1}), reflecting severe class imbalance. The weighted F1 is 0.590; the macro-averaged F1 (0.377) is degraded by the zero F1 on Grade I and limited F1 on Grade V.

\begin{figure}[!htbp]
    \centering
    \includegraphics[width=0.9\columnwidth]{figure9.pdf}
    \caption{Normalized Pfirrmann grading confusion matrix (A2, test set). Values are row-normalized proportions with counts in parentheses. Strong Grade II classification (87\%) with confusion primarily between adjacent grades. Grade I ($n=2$) is never correctly classified; Grade V ($n=8$) achieves 25\% recall.}
    \label{fig:confusion}
\end{figure}

\begin{figure}[!htbp]
    \centering
    \includegraphics[width=0.9\columnwidth]{figure10.pdf}
    \caption{Per-grade F1-score (blue) and sample size (orange) for Pfirrmann grading (A2, test set). The strong inverse relationship between F1 and class rarity highlights the need for class-balanced strategies for extreme grades.}
    \label{fig:grade_f1}
\end{figure}

\subsection{Disc Height Computation}
Complementing the grading analysis, we evaluate two additional clinical outputs: disc height regression and multi-label disorder detection. Disc height is a sagittal-plane measurement; ground-truth heights were derived from mid-sagittal T2-weighted slices using DICOM pixel spacing metadata (mean pixel spacing: $0.49 \pm 0.08$ mm) and the anterior-posterior midline protocol of Frobin \etal \cite{frobin1997precision}, then propagated as per-level labels to the corresponding axial images for multi-task training. Evaluation is performed on T2w axial images only ($n=234$, i.e. 78 patients $\times$ 3 disc levels), which avoids double-counting the same ground-truth measurement across T1w and T2w axial images of the same disc level. Disc height regression achieves MAE$=0.57$ mm with Pearson correlation $r=0.94$. Normal lumbar disc height ranges from approximately 7--12 mm \cite{guellil2022disc_metrics}, placing the 0.57 mm MAE at $\sim$3--8\% relative error, which is clinically informative for monitoring longitudinal change but insufficient for standalone diagnostic use. \Cref{fig:bland_altman} shows the Bland--Altman analysis: bias$={+}0.03$ mm with 95\% limits of agreement $[-1.67, {+}1.73]$ mm.

\begin{figure}[!htbp]
    \centering
    \includegraphics[width=\columnwidth]{figure11.pdf}
    \caption{Bland--Altman plot for disc height estimation (A2, T2w axial test images, $n=234$ unique disc levels; 78 patients $\times$ 3 disc levels). Ground-truth heights were measured from mid-sagittal T2-weighted images; predictions are from the corresponding T2w axial slice of each disc level. }
    \label{fig:bland_altman}
\end{figure}

% ------------------------------------------------------------------
\subsection{Disorder Detection}
\label{sec:disorder_results}
% ------------------------------------------------------------------

Multi-label disorder detection is evaluated at two levels: per-image (threshold $\theta=0.5$, $n=624$) and per-patient ($75^\text{th}$-percentile sigmoid aggregation across all views, $n=78$). Because radiologists render a single diagnostic impression per patient by synthesizing multiple views, the patient-level aggregation constitutes the more clinically relevant unit of inference. \Cref{tab:disorders} consolidates both evaluations alongside optimal per-disorder thresholds ($\theta^*$) identified by sweeping the validation set.

Common disorders are reliably detected at the baseline including thecal compression (per-patient F1 $= 0.830$), disc bulge (0.772) and protrusion (0.712) all exceed an F1 of 0.70 at the patient level. Furthermore, nerve compression improves significantly from a per-image F1 of 0.044 at $\theta=0.5$ to $0.721$ at $\theta^*=0.10$. This indicates that the sigmoid outputs cluster below the default 0.5 threshold but remain above the noise floor, representing a calibration artifact rather than a representational failure. Overall, patient-level aggregation raises the macro-averaged F1 from $0.237$ to $0.566$ ($+138\%$), reflecting the immense clinical benefit of pooling multi-view evidence prior to diagnosis.

Conversely, rare disorders (stenosis, spondylolisthesis and ligamentum flavum involvement; support $\leq 4$ patients) exhibit near-zero F1 scores under the default protocol. This failure mode is driven by extreme class imbalance rather than intrinsic feature limitation. By deploying a dedicated standalone ResNet-50 classifier equipped with focal loss ($\gamma=2$) and minority-class oversampling, meaningful detection capabilities are recovered for the rarest classes: stenosis F1 rises from $0.057$ to $0.286$ ($+400\%$) and ligamentum flavum from $0.091$ to $0.500$ ($+449\%$). These results confirm that per-disorder threshold calibration and patient-level aggregation are essential for clinical deployment, whereas extreme rare-class enhancement is best served by targeted loss-reweighting.

\begin{table*}[htbp]
\centering
\caption{%
Disorder detection summary (A2, test set). \textit{Per-image} columns evaluate individual images ($n=624$) using the default $\theta=0.5$ threshold; \textit{Per-patient} columns aggregate sigmoid scores at the $75^\text{th}$ percentile across all views ($n=78$ patients) using validation-calibrated thresholds ($\theta^*$). $\Delta_\theta$ denotes the F1 gain from threshold calibration alone. \textit{Enh. F1} reports the performance of a standalone classifier with focal loss ($\gamma=2$) and oversampling (* denotes targeted rare classes). \textbf{Bold} indicates F1 $\geq 0.70$.%
}
\label{tab:disorders}
\begin{tabular*}{\textwidth}{@{\extracolsep{\fill}} l c ccc c c ccc c @{}}
\toprule
& & \multicolumn{3}{c}{\textit{Per-image} ($\theta=0.5$)} 
& \multicolumn{2}{c}{\textit{Threshold calibration}} 
& \multicolumn{3}{c}{\textit{Per-patient}} 
& \\
\cmidrule{3-5} \cmidrule{6-7} \cmidrule{8-10}
Disorder & Support & P & R & F1 & $\theta^*$ & $\Delta_\theta$ & P & R & F1 & Enh. F1 \\
\midrule
Disc bulge & 392 & 0.66 & 0.83 & \textbf{0.73} & 0.10 & +0.04 & 0.63 & 1.00 & \textbf{0.77} & \textbf{0.77} \\
Thecal compression & 400 & 0.73 & 0.79 & \textbf{0.76} & 0.10 & +0.00 & 0.79 & 0.88 & \textbf{0.83} & \textbf{0.85} \\
Nerve compression & 352 & 0.64 & 0.54 & 0.58 & 0.10 & +0.68 & 0.61 & 0.61 & 0.61 & 0.63 \\
Disc protrusion & 200 & 0.68 & 0.32 & 0.44 & 0.10 & +0.05 & 0.62 & 0.84 & \textbf{0.71} & -- \\
Disc herniation & 112 & 0.43 & 0.21 & 0.28 & 0.50 & +0.00 & 0.37 & 0.71 & 0.49 & -- \\
Disc desiccation & 112 & 0.50 & 0.01 & 0.02 & 0.10 & +0.28 & 0.36 & 0.86 & 0.51 & -- \\
Stenosis* & 32 & 0.33 & 0.03 & 0.06 & 0.10 & +0.04 & 0.50 & 0.50 & 0.50 & 0.29 \\
Spondylolisthesis* & 16 & 1.00 & 0.06 & 0.12 & 0.50 & +0.00 & 0.50 & 1.00 & 0.67 & -- \\
Ligamentum flavum* & 16 & 0.17 & 0.06 & 0.09 & 0.50 & +0.00 & 0.00 & 0.00 & 0.00 & 0.50 \\
\midrule
\textbf{Macro Average} & -- & -- & -- & 0.34 & -- & +0.16 & -- & -- & \textbf{0.57} & 0.34 \\
\bottomrule
\end{tabular*}
\end{table*}


\subsection{Report Generation}
\label{sec:report_results}

Consolidating the preceding clinical outputs, \modelname{}'s final component is a retrieval-augmented template-based report generator. The Prompt Calibrator's learned vocabulary conditioning enables zero-shot report generation without additional training, while the retrieval mechanism enriches the context with similar cases from the training set to enhance clinical phrasing and content structure.

We comprehensively benchmark (\Cref{sec:report_results}) RATF against twelve foundation LLMs under a matched-patient protocol ($n=20$), spanning instruction-tuned general-purpose models (Phi-3.5-mini \cite{phi35mini2024}, Llama-3.2-3B-Instruct \cite{llama32_2024}, SmolLM2-1.7B-Instruct \cite{smollm2_2024} and Qwen2.5-1.5B-Instruct \cite{qwen25_2024}), domain-specialised medical language models (MedGemma-4B-IT \cite{medgemma2025}, BioMistral-7B \cite{biomistral2024}, Meditron3-8B \cite{sallinen2025llama3meditron}, BioGPT-Large \cite{luo2022biogpt} and Med-PaLM \cite{singhal2023medpalm2}) and radiology-specific generative models (LLaVA-Rad-7B \cite{zambranochaves2025llavarad}, RADAR \cite{sun2026radar} and MAIRA-2 \cite{bannur2024maira2}). 
% MIRA \cite{li2025mira} is excluded as its weights are currently unreleased. 

Each model received identical predicted measurements (grade, height, and disorders per level) to generate a full report. For instruction-tuned models, we utilized a structured chat-template prompt. For base language models, we employed a \emph{seeded-completion} strategy, pre-filling the standard clinical template to ensure well-formed outputs and isolate clinical reasoning from formatting adherence. Report quality was first assessed using standard NLG overlap metrics and contextual semantic similarity (BERTScore), computed against radiologist-authored reference notes (\Cref{tab:nlg_eval}). Lexical overlap metrics (BLEU-1 \cite{papineni2002bleu}, ROUGE-L \cite{lin2004rouge}) are uniformly low across all twelve baseline methods (BLEU-1 range: 0.024--0.103; ROUGE-L range: 0.039--0.150), reflecting the inherent vocabulary mismatch between free-text radiologist notes and algorithmically generated reports, a known and documented limitation widely discussed in the medical report generation literature.

Our proposed RATF substantially outperforms all baselines on every lexical metric (e.g., ROUGE-L $= 0.298$, $2.0\times$ the strongest baseline). By explicitly retrieving historical reports, RATF generates clinical prose that matches radiologist vocabulary patterns (e.g., ``compressing the thecal sac,'' ``encroaching upon exit nerve roots''). BERTScore \cite{zhang2020bertscore} further confirms this superiority (0.870), indicating that RATF reports are semantically closest to the reference texts. Furthermore, RATF's METEOR \cite{banerjee2005meteor} score ($0.384$) indicates promising lexical recall and synonym alignment, confirming that the retrieval-augmented template engine effectively preserves the morphological phrasing of clinical expert better than unconstrained generative baselines.

\begin{table}[htbp]
\centering
\caption{Comprehensive NLG evaluation of report generators ($n=20$ matched test patients). Reference reports are radiologist-authored clinical notes. Lexical metrics (BLEU, METEOR, ROUGE) measure surface-level overlap; BERTScore measures contextual semantic similarity; Factual F1 evaluates disorder-level clinical accuracy via synonym-aware keyword extraction. Best values are \textbf{bold}.}
\label{tab:nlg_eval}
\resizebox{\columnwidth}{!}{
\begin{tabular}{@{}l ccccc @{}}
\toprule
Method & BLEU-1 & METEOR & ROUGE-L & BERTScore & Factual F1 \\
\midrule
RATF (ours)        & \textbf{0.221} & \textbf{0.384} & \textbf{0.298} & \textbf{0.870} & \textbf{0.798} \\
Med-PaLM \cite{singhal2023medpalm2}          & 0.103 & 0.214 & 0.150 & 0.828 & 0.711 \\
Meditron3-8B \cite{sallinen2025llama3meditron}       & 0.060 & 0.199 & 0.106 & 0.813 & 0.699 \\
Llama-3.2-3B \cite{llama32_2024}      & 0.052 & 0.185 & 0.099 & 0.813 & 0.680 \\
BioMistral-7B \cite{biomistral2024}      & 0.097 & 0.237 & 0.134 & 0.827 & 0.655 \\
MedGemma-4B-IT \cite{medgemma2025}     & 0.039 & 0.156 & 0.086 & 0.794 & 0.652 \\
BioGPT-Large \cite{luo2022biogpt}       & 0.055 & 0.199 & 0.094 & 0.810 & 0.605 \\
LLaVA-Rad-7B \cite{zambranochaves2025llavarad}       & 0.092 & 0.207 & 0.143 & 0.830 & 0.602 \\
Qwen2.5-1.5B \cite{qwen25_2024}       & 0.078 & 0.238 & 0.138 & 0.829 & 0.602 \\
Phi-3.5-mini \cite{phi35mini2024}       & 0.069 & 0.194 & 0.112 & 0.819 & 0.587 \\
SmolLM2-1.7B \cite{smollm2_2024}       & 0.060 & 0.162 & 0.076 & 0.814 & 0.571 \\
MAIRA-2 \cite{bannur2024maira2}            & 0.084 & 0.137 & 0.112 & 0.831 & 0.220 \\
RADAR \cite{sun2026radar}             & 0.024 & 0.066 & 0.039 & 0.778 & 0.080 \\
\bottomrule
\end{tabular}
}
\end{table}

 
Beyond surface-level overlap, we quantified clinical content preservation using the Factual F1 metric (via synonym-aware keyword extraction), alongside computational efficiency (\Cref{tab:report_comparison}). RATF attains the highest factual consistency (F1 $= 0.798$) and the shortest mean generation time ($<0.1$ s per patient), reliably producing RSNA-compliant structured reports. However, it must be noted that RATF's clinical accuracy is inherently bounded by the upstream multi-task predictions; segmentation or classification errors propagate directly into the final text. 

Among the generative baselines, the benchmark establishes a clear performance hierarchy. Med-PaLM (F1 $= 0.711$), Phi-3.5-mini (F1 $= 0.705$), and Meditron3-8B (F1 $= 0.699$) proved the most robust. Meditron3-8B specifically demonstrates that medical-domain continued pretraining yields a meaningful advantage over general-purpose LLMs of comparable size. Furthermore, base models like BioMistral-7B and BioGPT-Large recovered to competitive factual performance only under the seeded-completion strategy, indicating that prompt structure, rather than intrinsic model reasoning, was their primary limiting factor. Conversely, domain mismatch remains the dominant failure mode for radiology-specific vision-language models. MAIRA-2 and RADAR, both trained exclusively on chest X-ray corpora, achieved the lowest factual F1 scores ($0.220$ and $0.080$), confirming that organ-specific training domains do not transfer to lumbar spine MRI reporting without targeted adaptation.

\begin{table}[htbp]
\centering
\caption{Report generation benchmark ($n=20$ matched test patients), ordered by factual F1 (ascending). Disorder precision (P), recall (R), and F1 are per-patient synonym-aware scores. Words: mean report length $\pm$ SD. $t_{\text{pt}}$: mean per-patient generation time. \textbf{Bold}: best value per column.}
\label{tab:report_comparison}
\resizebox{\columnwidth}{!}{
\begin{tabular}{@{}l ccccc @{}}
\toprule
Method & P$\uparrow$ & R$\uparrow$ & F1$\uparrow$ & Words & $t_{\text{pt}}$ \\
\midrule
RADAR \cite{sun2026radar}                 & 0.100 & 0.067 & 0.080 & $211{\pm}47$ & 103.47 s \\
MAIRA-2 \cite{bannur2024maira2}                & 0.425 & 0.158 & 0.220 & $80{\pm}59$  & 17.25 s \\
SmolLM2-1.7B-Instruct \cite{smollm2_2024}  & 0.613 & 0.650 & 0.596 & $124{\pm}7$  & 8.20 s \\
LLaVA-Rad-7B \cite{zambranochaves2025llavarad}           & 0.667 & 0.652 & 0.602 & $96{\pm}25$  & 15.76 s \\
BioGPT-Large \cite{luo2022biogpt}           & 0.505 & 0.867 & 0.605 & $168{\pm}65$ & 112.97 s \\
Qwen2.5-1.5B-Instruct \cite{qwen25_2024}  & 0.550 & 0.817 & 0.616 & $128{\pm}14$ & 9.63 s \\
MedGemma-4B-IT \cite{medgemma2025}         & 0.576 & 0.828 & 0.652 & $226{\pm}45$ & $\sim$200 s \\
BioMistral-7B \cite{biomistral2024}          & 0.607 & 0.832 & 0.655 & $95{\pm}9$   & 115.52 s \\
Llama-3.2-3B-Instruct \cite{llama32_2024}  & 0.648 & 0.807 & 0.673 & $203{\pm}15$ & 29.21 s \\
Meditron3-8B \cite{sallinen2025llama3meditron}           & 0.683 & 0.807 & 0.699 & $171{\pm}26$ & 23.35 s \\
Phi-3.5-mini \cite{phi35mini2024}           & 0.742 & 0.745 & 0.705 & $97{\pm}5$   & 77.95 s \\
Med-PaLM \cite{singhal2023medpalm2}               & 0.691 & 0.832 & 0.711 & $60{\pm}4$   & 22.42 s \\
\midrule
RATF (ours)            & \textbf{0.743} & \textbf{0.916} & \textbf{0.798} & $95{\pm}5$ & $<$\textbf{0.1 s} \\
\bottomrule
\end{tabular}
}
\end{table}

To evaluate the impact of the underlying embedding space, we replaced the general-domain all-MiniLM-L6-v2 encoder in the FAISS retrieval stage with MedCPT \cite{jin2023medcpt}, a biomedical text encoder contrastively pretrained on PubMed query-article pairs. On the matched test set, MedCPT improved retrieval F1 from $0.425$ to $0.441$ ($+3.8\%$ modest improvement). The higher-dimensional MedCPT embeddings (768-d vs. 384-d) encode richer biomedical semantics at a negligible increase in indexing time (2.2s vs. 0.7s for the full corpus), supporting its adoption as the default encoder in future iterations.

% \Cref{fig:report_comparison} presents a qualitative same-patient comparison of complete structured reports generated by \modelname{} and MedGemma-4B-IT for two representative test cases, illustrating the structural alignment of retrieval-augmented template filling with RSNA reporting conventions versus the free-form prose produced by end-to-end generative models.

% \begin{figure*}[htbp]
% \centering
% \includegraphics[width=\textwidth]{figure12a}
% \caption{Qualitative same-patient report comparison. Two representative test patients (Case A: progressive degeneration, Grade II$\to$III$\to$IV; Case B: multi-level Grade II with multiple disorders) are shown with their complete MRI inputs, comprising 6 axial images (T1 and T2, three disc levels each) and 2 sagittal images (T1 and T2), together with the structured RATF report generated by \modelname{} (left, \textcolor{NavyBlue}{navy}) and the MedGemma-4B-IT report (right, \textcolor{Violet}{purple}).
% Both reports are generated from the \emph{same} patient images without any cross-system data leakage.
% RATF integrates quantitative segmentation confidence (Dice), Pfirrmann grade, disc height and disorder predictions into a structured RSNA-compliant format; MedGemma generates natural-prose findings from image embeddings.}% All 20 matched-patient reports are available at \texttt{results/ratf\_reports\_matched/}.}
% \label{fig:report_comparison}
% \end{figure*}

The specific failure modes visualized in \Cref{fig:report_comparison} map directly onto the quantitative performance gaps reported in \Cref{tab:report_comparison}. For instance, L3-L4, L4-L5 IVD level hallucination clearly visible in MedGemma-4B-IT's output reveal its dependency on prompt completeness. \modelname{} circumvents this issue by strictly grounding every report field in upstream deterministic predictions rather than relying on open-ended generation. Similarly, BioGPT-Large's visually replicated IVD clinical findings and contradictory impression grades are known manifestations of the repetition artifacts documented for GPT-2 architectures \cite{luo2022biogpt}. Finally, the visible domain failure illustrated by RADAR and MAIRA-2, generating anatomically irrelevant thoracic and planar chest-radiograph text for a lumbar MRI highlights a systemic limitation of anatomy-specific vision-language models. High performance on a source domain provides no transferable clinical vocabulary for a distinct target; a critical gap that \modelname{} explicitly resolves by indexing its retrieval-augmented engine exclusively against target-specific clinical notes.

\begin{figure*}[htbp]
    \centering
    \includegraphics[width=\textwidth]{figure12a}
    \caption{Qualitative comparison of Report Generation across five models. \modelname{} (orange) generates a structured, RSNA-compliant report with per-level Pfirrmann grades, disc height estimates, disorder findings and per-image segmentation confidence (\textit{Dice}), demonstrating differentiated level-specific assessment. MedGemma-4B-IT (yellow) produces fluent prose but over simplifies the individual assessment. BioGPT-Large (green) replicates identical clinical notes for IVD assessment and outputs contradictory grade labels in the Impression (Grade 0) versus Findings (Grades 2--4) sections. RADAR (blue) and MAIRA-2 (pink) exhibit complete domain failure, generating thoracic and chest-radiograph content irrelevant to lumbar spine pathology, consistent with their training on MIMIC-CXR/CheXpert-Plus datasets.}
    \label{fig:report_comparison}
\end{figure*}
% ──────────────────────────────────────────────────────────────────
\subsection{Cross-Dataset Validation}
\label{sec:external_validation}

To assess generalizability beyond LSMA-PQR (internal dataset), we evaluated the \modelname{} architecture \emph{zero-shot} (without any fine-tuning or domain adaptation) on three external MRI datasets. To ensure transferability, all external datasets were processed with a standardized pipeline for extraction of mid-sagittal slices, spatial resampling to $384{\times}384$, matching the training protocols and the clinical reference plane for spine MRI diagnosis. Crucially, zero-shot inference was executed utilizing \emph{per-class text prompting} and \emph{plane-aware head selection}. For each expected foreground class, inference was conditioned on its canonical prompt from the vocabulary bank, ensuring the Prompt Calibrator operated identically to its training regime. Furthermore, anatomical head routing was strictly enforced: the sagittal head ($\hat{S}_{\text{sag}}$, isolating four classes including VB and IVD\textsubscript{sag}) was activated for all sagittal MRI series, while the axial head ($\hat{S}_{\text{ax}}$, isolating three classes) was restricted to axial T2 series. Pfirrmann grading was constrained exclusively to mid-sagittal slices, ensuring clinical fidelity to standard disc signal and morphological assessment protocols. 

The aggregated zero-shot performance across all external datasets is summarized in \Cref{tab:external_summary}. To ensure fair comparative analysis, several dataset-specific evaluation constraints were enforced. For the SPIDER dataset \cite{van2024spider}, evaluation was strictly restricted to the 210 mid-sagittal slices to match our training distribution and standard clinical assessment planes, with the disorder F1-score specifically reflecting disc bulge detection. For RSNA 2024 \cite{rsna_lumbar_2026}, performance was evaluated across 500 individual disc levels utilizing a Pfirrmann-to-stenosis proxy mapping, where the reported F1-score isolates axial spinal canal stenosis. On SpineWeb \cite{spineweb2012} and the SpineParseNet Dataset \cite{spineparsenet2021}, analysis was constrained exclusively to vertebral body (VB) segmentation. 

The SPIDER cohort \cite{van2024spider} comprises 210 T2-weighted sagittal MRI volumes with dense voxel-level annotations. We mapped the provided labels to our sagittal taxonomy (VB, IVD, SC) and evaluated the model strictly on the mid-sagittal slices ($n=210$) to preserve clinical alignment. As detailed in \Cref{tab:external_perclass}, the spinal canal (SC\textsubscript{sag}) exhibited the most robust cross-dataset transfer (Dice = 0.705), followed by vertebral bodies (VB Dice = 0.643). The mean foreground Dice across the cohort settled at 0.616. Notably, clinical-task predictions transferred exceptionally well: zero-shot Pfirrmann grading achieved a 75.2\% within-$\pm 1$ accuracy across 1,461 disc levels, and disc bulge detection reached an F1-score of 0.848 (with 94.3\% precision) across 195 positive cases.

To assess multi-institutional robustness, we evaluated 100 studies (500 disc levels) from the RSNA 2024 dataset \cite{rsna_lumbar_2026}. Because this dataset lacks Pfirrmann labels but provides spinal canal stenosis severity (Normal/Mild, Moderate, Severe), we employed an ordinal proxy mapping (Grades I--II $\to$ Normal/Mild; Grade III $\to$ Moderate; Grades IV--V $\to$ Severe). While Pfirrmann grading directly assesses disc degeneration rather than canal narrowing (making this an imperfect proxy), the model demonstrated strong ordinal-scale consistency, achieving 39.4\% exact accuracy and 95.0\% within-$\pm 1$ accuracy. Concurrently, the axial stenosis detection head recorded a binary F1-score of 0.23 (Precision = 0.55), reflecting the difficulty of zero-shot transfer for severe minority-class pathologies.

Generalization for vertebral body (VB) segmentation was further validated on two highly heterogeneous public benchmarks. The SpineParseNet Dataset \cite{spineparsenet2021}, which features exceptionally thick slices (4.4 mm) and high in-plane resolution (0.345 mm spacing), yielded a zero-shot VB Dice of 0.367 across 633 foreground-containing slices. Conversely, the SpineWeb dataset \cite{spineweb2012}, spanning diverse T1, T2, and TIRM sequences across 209 slices, yielded a highly competitive VB Dice of $0.668 \pm 0.157$ (IoU = 0.519). This confirms that the sagittal segmentation head maintains robust feature extraction capabilities across highly variable acquisition physics and multi-institutional protocols.

% {\color{BrickRed}
% Finally, to probe the modality boundaries of the learned representations, the framework was applied to an dataset of 50 anterior-posterior scoliosis radiographs. The model predicted negligible foreground area (<2\% pixel density for all classes). This catastrophic failure confirms that while the framework is robust to MRI sequence variations, the underlying latent space does not natively bridge distinct imaging modalities (MRI to X-ray) without explicit cross-modal domain adaptation, a phenomenon consistent with prevailing literature on medical foundation models \cite{ma2024medsam}.}

% Finally, the internal scoliosis X-ray cohort (comprising anterior-posterior radiographs collected under the same ethics protocol as LSMA-PQR) lacks ground-truth segmentation annotations, precluding the computation of a Dice score; consequently, performance on this modality is reported via per-class predicted area fractions in subsequent analyses.

\begin{table}[htbp]
\centering
\caption{Cross-dataset zero-shot evaluation of the \modelname{} framework. Evaluation relies strictly on per-class text prompting and plane-aware head selection without downstream fine-tuning. $\bar{D}$: Mean foreground Dice; W-1: Within-$\pm 1$ grade accuracy; Dis-F1: Primary disorder detection F1-score.}
\label{tab:external_summary}
\resizebox{\columnwidth}{!}{%
\begin{tabular}{@{}l c c c cccc @{}}
\toprule
Dataset & Mod. & $n$ & Classes & $\bar{D}\uparrow$ & Accuracy$\uparrow$ & W-1$\uparrow$ & Dis-F1$\uparrow$ \\
\midrule
LSMA-PQR (Ours)                     & MRI   & 624 & 7 (all) & 0.861 & 61.8\% & 92.3\% & 0.634 \\
SPIDER \cite{van2024spider}         & MRI   & 210 & 3 (sag) & 0.616 & 29.6\% & 75.2\% & 0.848 \\
RSNA 2024 \cite{rsna_lumbar_2026}   & MRI   & 500 & -    & -  & 39.4\% & 95.0\% & 0.230 \\
SpineWeb \cite{spineweb2012}        & MRI   & 209 & 1 (VB)  & 0.668 & -    & -    & - \\
SpineParseNet Dataset \cite{spineparsenet2021} & MRI & 633 & 1 (VB)  & 0.366 & -   & -    & - \\
% Internal Scoliosis Cohort           & X-ray & 50  & -    & n/a   & -    & -    & - \\
\bottomrule
\end{tabular}%
}
\end{table}



To further quantify \modelname{}'s structural generalization, \Cref{tab:external_perclass} compares per-class segmentation performance between the LSMA-PQR test set and the zero-shot SPIDER dataset \cite{van2024spider}. SPIDER evaluation is constrained to the central mid-sagittal slice per volume and Dice coefficients are computed on merged class-union binary masks. The spinal canal (SC\textsubscript{sag}) demonstrates the highest resilience to cross-domain distribution shifts, maintaining a zero-shot Dice score of $0.705$. Vertebral body (VB) segmentation achieves moderate transferability ($0.643 \pm 0.226$), while the intervertebral disc (IVD\textsubscript{sag}) exhibits the steepest performance degradation ($0.500 \pm 0.168$). Furthermore, the substantially wider standard deviations observed across all SPIDER classes highlight the inherent variance of zero-shot inference on external clinical acquisition protocols, contrasting sharply with the highly stable in-distribution baseline ($\sigma < 0.05$). Despite the domain shift, the \modelname{} retains a mean foreground Dice of $0.616$ on the unseen cohort without requiring any dataset-specific fine-tuning. In the \Cref{tab:external_perclass}, Dice scores are computed on merged class-union binary masks per slice (\ie, all vertebral bodies in a slice are merged into a single mask). The slice count is identical across all classes within each dataset. $^{\dagger}$Mean foreground for SPIDER is averaged over the three available classes, as Sacrum annotations are not provided in the target dataset



\begin{table}[htbp]
\centering
\caption{Per-class segmentation Dice ($D \pm \sigma$) comparing \modelname{}'s in-distribution performance against zero-shot transfer to the SPIDER dataset \cite{van2024spider}.}
\label{tab:external_perclass}
\resizebox{\columnwidth}{!}{%
\begin{tabular}{@{}l cc@{}}
\toprule
& \textbf{Internal Dataset} & \textbf{External Dataset (Zero-Shot)} \\
Anatomical Class & LSMA-PQR ($n=156$) & SPIDER ($n=210$) \\
\midrule
Vertebral Body (VB)                          & $0.899 \pm 0.034$ & $0.643 \pm 0.226$ \\
Intervertebral Disc (IVD\textsubscript{sag}) & $0.824 \pm 0.048$ & $0.500 \pm 0.168$ \\
Spinal Canal (SC\textsubscript{sag})         & $0.886 \pm 0.036$ & $0.705 \pm 0.272$ \\
Sacrum                                       & $0.892 \pm 0.043$ & --- \\
\midrule
\textbf{Mean Foreground}                     & $0.875$           & $0.616^{\dagger}$ \\
\bottomrule
\end{tabular}
}
\end{table}


\subsection{Qualitative Results}
After establishing the quantitative performance of \modelname{} across all tasks, we present representative qualitative results to illustrate the clinical relevance and anatomical coherence of the model's predictions.

\Cref{fig:external_qual} presents representative cross-dataset predictions on SPIDER \cite{van2024spider} and SpineWeb \cite{spineweb2012} volumes, showing that the ROI-cropped, per-class-prompted evaluation yields anatomically coherent segmentation masks that capture the major structural boundaries, despite no exposure to these datasets during training.


\begin{figure}[!htbp]
    \centering
    \includegraphics[width=\columnwidth,height=0.62\textheight,keepaspectratio]{figure13.png}
    \caption{Qualitative cross-dataset segmentation results (zero-shot, no fine-tuning, mid-sagittal view). Rows: representative SPIDER T2 mid-sagittal slices and SpineWeb Datasets sagittal slices. Columns: (a) input image resized to $384{\times}384$, (b) ground-truth annotation, (c) \modelname{} prediction using the sagittal head (plane\_id$=1$).}% SpineWeb visualisations ($D_\text{VB}=0.668$) and Scoliosis X-ray results ($\bar{D} < 0.02$) are provided in supplementary figures.}
    \label{fig:external_qual}
\end{figure}



\Cref{fig:qualitative} shows representative segmentation results on both axial and sagittal planes.
The model accurately delineates major structures (IVD, VB) with smooth boundaries, while smaller structures (SC) remain challenging.

\begin{figure}[!htbp]
    \centering
    \includegraphics[width=\columnwidth,height=0.72\textheight,keepaspectratio]{figure13_qual.pdf}
    \caption{Qualitative segmentation results on validation samples (\texttt{seed=42}, no cherry-picking; shown instead of test samples to preserve the held-out test set for quantitative evaluation only). Columns: (a) input MRI, (b) ground-truth multi-class mask overlay, (c) predicted prompted mask with Dice score, (d) ground-truth prompted mask. The model accurately segments major structures on both axial and sagittal planes.}
    \label{fig:qualitative}
\end{figure}

\begin{figure*}[!htbp]
    \centering
    \includegraphics[width=\textwidth,height=0.62\textheight,keepaspectratio]{figure14_failure.pdf}
    \caption{Failure case gallery (A2, test set). Representative failure modes: (a) axial spinal canal under-segmentation due to extreme minority pixel ratio ($<$0.25\%), (b) boundary confusion between IVD and adjacent VB in sagittal views, (c) false-negative disorder prediction for rare pathologies (stenosis, spondylolisthesis), (d) Pfirrmann grade boundary errors (Grade III predicted as Grade II). These failure modes correlate with the quantitative limitations identified in \Cref{sec:grading_results} and \Cref{tab:disorders}.}
    \label{fig:failure_gallery}
\end{figure*}

% ══════════════════════════════════════════════════════════════════
\section{Discussion}
\label{sec:discussion}

The experimental results demonstrate that text-prompted, annotation-free segmentation of lumbar spine MRI is achievable within a unified multi-task architecture at a computational cost compatible with standard clinical hardware. The sections below discuss the architectural contributions, clinical output quality and cross-dataset performance, followed by a review of current limitations and directions for future development.

A core contribution of \modelname{} is the elimination of spatial annotations during inference. Existing approaches require ground-truth spatial prompts, whereas \modelname{} achieves highly competitive segmentation ($\bar{D}=0.861$) while relying entirely on clinical text prompts. This capability is driven by the Prompt Calibrator, which maps closed-vocabulary clinical prompts to channel-wise feature modulation.  

The ablation studies (\Cref{tab:ablation}) quantitatively reinforce the necessity of this prompt-conditioned gating. Applying low-rank adaptation alone without the Prompt Calibrator yields no improvement in prompted segmentation, demonstrating that standard feature adaptation is insufficient for text-prompted spine MRI interpretation. Integrating the Prompt Calibrator yields a substantial Dice improvement of $+0.486$. Interestingly, the isolated low-rank adaptation model achieves the highest exact grade accuracy ($0.652$) despite demonstrating near-random segmentation ($\bar{D}=0.375$) performance. It is because the CORN grade head operates on mean-pooled encoder features independently of the segmentation decoder, LoRA adaptation can improve grade-relevant feature representations without producing useful segmentation outputs. In the fully integrated model, the segmentation and grading losses share the same feature space, redistributing capacity across tasks. While this introduces mild gradient competition, the marginal decrease of $0.034$ in exact grading accuracy is heavily outweighed by the impressive gains of $+0.486$ in spatial segmentation, validating the multi-task design. Furthermore, tuning the boundary loss weight reveals a consistent multi-task trade-off. Increasing the boundary penalty marginally improves the delineation of thin structures like the axial spinal canal but simultaneously degrades Pfirrmann grading accuracy and height estimation. A conservative boundary weight ultimately provides the most stable equilibrium across all clinical tasks.

To empirically validate the representational impact of the multi-conditional Prompt Calibrator, we evaluated the topological structure of the decoder's feature manifold using t-SNE \cite{vandermaaten2008tsne} and UMAP projections (\Cref{fig:feature_space}). In the absence of text-conditioned gating (Panels A and C), the feature space collapses into a homogenous, entangled distribution. This spatial ambiguity is mathematically confirmed by a negative Silhouette coefficient ($sil = -0.030$), indicating that intra-class variance actively exceeds inter-class distance. Consequently, downstream task heads operating on this uncalibrated baseline are functionally ineffective to distinct anatomical boundaries, forced to rely on weak spatial priors rather than semantic understanding.

Conversely, the application of the A3 prompt gate (Panels B and D) triggers a profound topological disentanglement. The shared feature space explicitly splits into cleanly separable, structure-specific sub-manifolds. Quantitatively, the A3 gate elevates the Silhouette score to $+0.863$ for t-SNE and $+0.921$ for UMAP. Furthermore, the calibrated manifold successfully isolates orthogonal acquisition planes, clearly separating axial instances (\eg, IVD, SC) from their sagittal counterparts (\eg, IVD\_sag, SC\_sag). This confirms that the Prompt Calibrator successfully encodes the complete multi-conditional semantic space ($\Pi \times \mathcal{A}$), ensuring that the downstream grading and segmentation heads receive highly orthogonalized, non-overlapping feature sets.

\begin{figure}[!htbp]
\centering
\includegraphics[width=\columnwidth]{figure16_tsne.pdf}
\caption{t-SNE \cite{vandermaaten2008tsne} and UMAP feature-space visualizations of \modelname{} backbone representations, comparing ungated features (A,C) against Prompt-Calibrator-gated features (B,D) for all seven anatomical classes (test set with $n=624$ samples). Without gating, all classes collapse into two indistinct blobs stratified only by imaging plane (axial vs. sagittal), yielding a silhouette coefficient of $-0.030$ in both projections. With gating (A3), each class forms a tight, well-separated cluster: silhouette rises to $+0.863$ under t-SNE and $+0.921$ under UMAP.}
\label{fig:feature_space}
\end{figure}


\Cref{fig:prompt_gate} examines how the Prompt Calibrator in \modelname{} converts a text condition into selective feature modulation. Panel A reveals that the 768-dimensional DINOv2 feature space is partitioned into two largely non-overlapping functional subsets: axial-active channels, driven by axial anatomy conditions (IVD, PE, SC), and sagittal-active channels, driven by sagittal conditions (VB, IVD\textsubscript{sag}, Sacrum, SC\textsubscript{sag}). Axial conditions activate 12-15 channels and suppress sagittal-active channels to near zero; sagittal conditions show the complementary pattern, activating 16-19 channels, with the higher count reflecting the greater structural diversity present in mid-sagittal views. PanelB quantifies the overlap between gate vectors via cosine similarity: within-plane pairs are consistently more similar (axial-axial: 0.56-0.74; sagittal-sagittal: 0.69-0.76) than cross-plane pairs (0.27-0.60). The minimum cross-plane similarity of 0.27 occurs between Axial PE and Sag IVD\textsubscript{sag} posterior elements appear in axial slices as ring-shaped cortical structures, whereas the sagittal IVD appears as a horizontally banded ovoid, making their gate vectors nearly orthogonal. The highest cross-plane similarity of 0.60 occurs between Axial IVD and Sag IVD\textsubscript{sag}, consistent with both conditions representing the same disc tissue viewed from orthogonal planes.

\begin{figure*}[!htbp]
    \centering
    \includegraphics[width=\textwidth]{figure15_gate_cosine.pdf}
    \caption{Analysis of the Prompt Calibrator gating mechanism (\modelname{}).
    \textbf{(A)} Sigmoid gate values $\sigma(\mathbf{W}_g \mathbf{c})$ for the top 32 channels by across-condition variance. Bar charts report active-condition counts per channel (bottom) and active-channel counts per condition (right). \textbf{(B)} Inter-condition cosine similarity between learned gate vectors.}
    \label{fig:prompt_gate}
\end{figure*}

Evidence from \Cref{fig:prompt_gate} analyzed together with the silhouette coefficient improvement from $-0.030$ (no gating) to $+0.863$ under t-SNE and $+0.921$ under UMAP (\Cref{fig:tsne_umap_gating}), confirm that the Prompt Calibrator performs anatomy and plane-discriminative channel selection rather than uniform amplification.

% \begin{figure*}[!htbp]
%     \centering
%     \includegraphics[width=\textwidth]{figure15_gate_cosine.pdf}
%     \caption{Prompt Calibrator gate analysis. Per-channel sigmoid gate values $\sigma(\mathbf{W}_g \mathbf{c})$ for each anatomy$\times$plane condition (top 32 channels by variance). The heatmap reveals structure-specific activation patterns: axial classes (IVD, PE, SC) activate distinct channel subsets from sagittal classes (VB, IVD\textsubscript{sag}, Sacrum, SC\textsubscript{sag}), confirming that the Prompt Calibrator learns anatomy-discriminative feature modulation rather than uniform amplification.}
%     \label{fig:prompt_gate}
% \end{figure*}



From a clinical deployment perspective, \modelname{} operates efficiently under strict hardware constraints (\Cref{tab:baselines}). Utilizing only 2.73 million trainable parameters ($3.08\%$ of total) and 0.63 GB of peak VRAM, our model executes the entire multi-task pipeline involving segmentation masks, Pfirrmann grades, disc height estimates and disorder probabilities in a single forward pass consuming 39 milliseconds per image on standard clinical workstations. For the subsequent reporting phase, the RATF module achieves the strongest factual consistency among all benchmarked generative models (\Cref{tab:report_comparison}). By retrieving historical templates, the module generates clinical prose that faithfully matches radiologist vocabulary patterns, avoiding the severe vocabulary mismatch and hallucinated formatting artifacts commonly observed in unconstrained foundational language models. Additionally, replacing the general-domain all-MiniLM-L6-v2 sentence encoder with MedCPT \cite{jin2023medcpt} improves retrieval F1 from $0.425$ to $0.441$ (+3.8\%), confirming that a biomedical-domain retrieval encoder benefits precision for spine pathology terminology.

Considered jointly, the external validation results (\Cref{tab:external_summary}) delineate two core structural properties governing the zero-shot generalization capabilities of \modelname{}. View alignment strictly dictates the efficacy of segmentation transfer. Restricting the SPIDER \cite{van2024spider} evaluation exclusively to the mid-sagittal slice, which mirrors the anatomical plane established in the training distribution, yields a robust mean foreground Dice of $0.616$. This spatial constraint correctly isolates the clinically informative central region and avoids a naive evaluation that inherently conflates target anatomy with uninterpretable peripheral planes. Moreover, ordinal clinical task predictions demonstrate substantially greater cross-domain resilience than dense pixel-level segmentation. Ordinal Pfirrmann grading retains high clinical utility across varying institutions, achieving 75.2\% within 1-grade accuracy on the SPIDER dataset and 95.0\% on the RSNA dataset \cite{rsna_lumbar_2026}. Conversely, spatial segmentation fidelity degrades proportionally to the divergence in scanner acquisition protocols. While boundary delineation remains moderate on SPIDER and SpineWeb \cite{spineweb2012}, performance drops significantly on the Spine DataSet NEW \cite{spineparsenet2021}. This pronounced degradation directly correlates with the atypical 880$\times$880 in-plane resolution, representing a structural domain shift from the native 384$\times$384 training protocol of \modelname{}.

% Limitations

Several limitations within the current framework merit consideration. First, while the \modelname{} achieves high computational efficiency, dedicated single-task segmentation baselines (such as nnU-Net ($0.944$) and SwinUNETR ($0.937$)) achieve higher raw Dice coefficients. This is expected as \modelname{} optimizes five output heads with $11.5{\times}$ fewer trainable parameters than the single-task alternatives, thereby trading absolute segmentation precision for increased clinical breadth and computational efficiency. Additionally, the clinical accuracy of the generated text reports is inherently bounded by the upstream multi-task predictions. Any misclassification in grading or disorder detection propagates directly into the final clinical note. 

Second, \modelname{} struggles with extreme class imbalances. Rare clinical findings, such as specific manifestations of stenosis and ligamentum flavum hypertrophy, exhibit low detection recall. Similarly, Pfirrmann Grades I and V contain insufficient test samples to reliably estimate per-class accuracy. While targeted focal loss reweighting and coarse-to-fine anatomical zooming partially recover performance for minority classes, these strategies have yet to be fully integrated into a unified training objective. 

Finally, the cross-dataset segmentation gap highlights a sensitivity to institutional acquisition protocols. Although per-class prompting and plane-aware head selection substantially reduce the evaluation gap on external datasets, the residual performance drop indicates that \modelname{} cannot yet provide deployment-grade segmentation across entirely unseen hospitals without localized calibration. Furthermore, \modelname{} struggles to diagnose scoliosis cases in cross-modal screening using coronal radiographs, confirming the requirement for domain adaptation across imaging modalities. 

% Future work

Building upon these findings, several extensions are identified as priorities. Reducing the residual cross-dataset segmentation gap will require training data drawn from multiple sites and scanner protocols; test-time augmentation or site-specific adapter modules may provide incremental improvements as additional annotated data from diverse acquisition settings becomes available. Extension to CT and radiograph modalities will require modality-specific adaptation, as the current architecture is trained exclusively on MRI; recent work on spine radiograph segmentation provides a relevant starting point \cite{spinefm2025}. Expanding the anatomical scope is another primary objective. Following the suggested domain adaptation protocols (Appendix \ref{app:generalization}), the modular architecture can be systematically extended (\Cref{alg:adaptation}) to the cervical spine, knee and cardiac MRI, provided sufficient localized annotations are curated. Moreover, extending the representations to process 3D volumetric context through inter-slice attention mechanisms could resolve boundary discontinuities. Finally, the report generation pipeline can be further improved by adopting biomedical-specific retrieval encoders and rigorous clinical validation through direct radiologist review of generated reports. 

\section{Conclusion}
\label{sec:conclusion}

We presented \modelname{}, a novel unified framework to address a practical gap in automated lumbar spine MRI analysis: the absence of a single, parameter-efficient model capable of jointly producing segmentation masks, Pfirrmann grades, disc height estimates, disorder probabilities and a structured radiology report from a clinical text prompt alone, without requiring spatial annotations at inference. The prompt calibrator's multi-conditional feature gating ensures superior formal expressiveness over single-conditional variants, while the Conditional Ordinal Regression loss intrinsically enforces monotonic grade probabilities. Together, these mechanisms establish a principled mathematical foundation for the robust joint performance observed across classification, regression and language generation tasks. Quantitatively, \modelname{} achieves mean Dice 0.788 across seven anatomical classes on the test set, Pfirrmann grade accuracy of 61.8\% exact and 92.3\% within one grade, disc height mean absolute error of 0.57 mm and patient-level disorder macro-F1 of 0.57, all within a single 39 ms forward pass requiring 2.73 M trainable parameters and 0.63 GB peak memory. The retrieval-augmented report generator reaches the highest F1 of 0.798 among thirteen benchmarked systems. Validation on external datasets demonstrates that clincial predictions transfer robustly across institutions. 

Overall our work demonstrates that the computational cost of full multi-task lumbar spine MRI analysis is not a barrier to clinical deployment. A single lightweight inference pass, driven by standard text terminology, replaces five isolated annotation-dependent workflows. Integration of these capabilities into a unified parameter-efficient architecture provides the clinical AI community a viable pathway toward deployment-grade automated spinal diagnostics at scale.


% ══════════════════════════════════════════════════════════════════
% APPENDICES
% ══════════════════════════════════════════════════════════════════
\appendices

\section{Domain Adaptation Protocol and Roadmap}
\label{app:generalization}

The modular architecture of \modelname{} supports principled adaptation to new anatomical domains while retaining its zero-shot prompting capabilities. Transferring the architecture to an unseen domain requires the reconfiguration of six specific components:

\begin{enumerate}[label=(\roman*), leftmargin=*]
    \item Replacement of Spinal terms (Vocabulary ($\mathcal{V}$)) with domain-specific anatomy.
    \item Updation of the acquisition geometry and resizing of the plane embedding (imaging planes ($\Pi$)).
    \item Substitution of clinical modules (task heads).
    \item Recomputation of $w_c$ based on the new dataset's pixel distribution (class weighting).
    \item Mapping of canonical text synonyms to the new clinical taxonomy (prompt bank).
    \item Replacement of the Spine Domain Adapter with a target-domain equivalent for compatibility of region-specific radiological features into the backbone.
\end{enumerate}

Crucially, the foundational feature extractors of \modelname{}, including the frozen DINOv2 backbone, the general LoRA adapters, and the RATF report engine, transfer directly without structural modification. The adaptation procedure (\Cref{alg:adaptation}) employs a two-phase parameter-efficient fine-tuning (PEFT) strategy designed to minimize data annotation burdens. 

\begin{algorithm}[htbp]
\caption{\modelname{} Domain Adaptation Protocol}
\label{alg:adaptation}
\begin{algorithmic}[1]
\Require Dataset $\mathcal{D}_{\text{new}}$, vocabulary $\mathcal{V}_{\text{new}}$, planes $\Pi_{\text{new}}$, tasks $\mathcal{T}$, checkpoint $\theta_{\text{source}}$.
\State \textbf{// Phase 0: Configuration}
\State Initialize $\mathcal{B}_{\text{new}}$, $\text{Emb}_{\text{anat}}$, $\text{Emb}_{\text{view}}$ for target domain.
\State Adapt task head dimensions from $\mathcal{T}$ and recompute $w_c$.
\State \textbf{// Phase 1: Head Warm-Up (5 epochs)}
\State \textbf{Freeze}: DINOv2 backbone, LoRA adapters, Domain Adapter, Decoder trunk.
\State \textbf{Train}: Emb$_{\text{anat}}$, Emb$_{\text{view}}$, Task Heads, Prompt Calibrator MLP (lr $= 10^{-3}$).
\State \textbf{// Phase 2: Full PEFT Fine-Tuning (up to 50 epochs)}
\State \textbf{Unfreeze}: LoRA adapters, Domain Adapter, Prompt Calibrator.
\State \textbf{Train} via AdamW (lr $= 5 \times 10^{-5}$) with cosine annealing and early stopping.
\State \Return Adapted model $\theta_{\text{new}}$
\end{algorithmic}
\end{algorithm}

Adapting the framework to structurally divergent domains, such as brain or cardiac MRI, requires redefining the target vocabulary ($|\mathcal{V}|$), spatial classes ($C$) and ordinal levels ($K$) alongside moderate Phase-2 fine-tuning ($n \approx 1,200\text{--}2,500$); additionally, cross-modality shifts to CT or planar radiographs necessitate full phase-2 fine tuning due to extensive distributional divergence.

% ══════════════════════════════════════════════════════════════════
% BIBLIOGRAPHY
% ══════════════════════════════════════════════════════════════════
\bibliographystyle{IEEEtran}
\bibliography{references}

\end{document}
